\documentclass[%
 reprint,
superscriptaddress,
 amsmath,amssymb,
 aps,
prb,
]{revtex4-2}

\usepackage{graphicx}% Include figure files
\usepackage{dcolumn}% Align table columns on decimal point
\usepackage{bm}% bold math
\usepackage{tikz}

\usepackage[utf8]{inputenc}
\usepackage[T1]{fontenc}
\usepackage{mathptmx}
\usepackage{siunitx}
\usepackage{amsmath}
\usepackage{array}
\usepackage{booktabs}
\usepackage[normalem]{ulem}
\usepackage{xcolor}
\usepackage[ colorlinks,
    linkcolor={blue},
    citecolor={blue},
    urlcolor={blue}]{hyperref}
\usepackage{braket}

%%%%%%%%%% Macros of the equations document (gamma_only_rirs_gw_equations.tex)
\newcommand{\br}{\mathbf{r}}
\newcommand{\bR}{\mathbf{R}}
\newcommand{\bq}{\mathbf{q}}
\newcommand{\bG}{\mathbf{G}}
\newcommand{\bk}{\mathbf{k}}
\newcommand{\qhat}{\hat{\mathbf{q}}}
\newcommand{\dd}{\mathrm{d}}
\newcommand{\ii}{\mathrm{i}}
\newcommand{\ee}{\mathrm{e}}
\newcommand{\MIC}{\mathrm{MIC}}
\newcommand{\Tr}{\mathrm{T}}
\newcommand{\head}{\mathrm{h}}
\newcommand{\wing}{\mathrm{w}}
\newcommand{\lat}{\mathrm{lat}}
\newcommand{\reg}{\mathrm{reg}}
\newcommand{\occ}{\mathrm{occ}}
\newcommand{\vir}{\mathrm{vir}}
\newcommand{\dens}{\mathrm{dens}}
\newcommand{\ext}{\mathrm{ext}}
\newcommand{\tc}{\mathrm{tc}}
%%%%%%%%%% Article-only macros
\newcommand{\hd}{\mathrm{hd}}
\newcommand{\Sov}{\mathsf{S}}% AO overlap matrix (distinct from the head tensor S_ab)
\newcommand{\cpk}{\textsc{cp2k}}

\begin{document}

\title{$\Gamma$-point head-and-wing correction for $\Gamma$-only real-space RI $GW$\\
of two- and three-dimensional periodic systems}

\author{Ritaj Tyagi}
\affiliation{University of Regensburg, Germany}

\date{November 2026}

\begin{abstract}
Large-cell $GW$ calculations of defects, interfaces and moir\'e superlattices are naturally
done at the $\Gamma$ point of a supercell, but the screened interaction $W$ at the single
sampled wavevector $\bq=0$ is ill defined: the real-space lattice sum of the bare Coulomb
operator between charged auxiliary functions diverges there. We present a scheme that makes
$\Gamma$-only $GW$ with real-space resolution of the identity (RI-RS) well defined for
two-dimensional slabs and three-dimensional bulk cells. The divergence is confined to a single
direction of the auxiliary basis, the charge channel. In two dimensions it is removed from
$V(0)$ as an exactly computable rank-one constant; in three dimensions $V(0)$ is taken as an
Ewald sum with the $\bG=0$ term omitted. The charge channel is then carried analytically by one
additional basis function, the normalized constant, whose screened head and wings are obtained
in closed form from velocity-gauge transition dipoles and the Gaussian dielectric body and are
integrated over the whole small Brillouin zone of the supercell. Nothing is subtracted, because
the single sampled point carries no head once the constant has been removed. In two dimensions
the correction is applied to $W$ in the auxiliary basis; in three dimensions, where the head is
eV-sized, it is applied to the correlation self-energy in the orbital basis through a single
orbital-pair kernel that contains the head, the head-dipole and the wing terms. We give every
working equation in the order of evaluation, describe the implementation in \cpk{}, and discuss
why the earlier $k$-mesh large-cell scheme is not a valid reference for two-dimensional systems.
\end{abstract}

\maketitle

%======================================================================================
\section{Introduction}
\label{sec:intro}
%======================================================================================

Many-body perturbation theory in the $GW$ approximation is the standard route to quasiparticle
band structures of solids, surfaces and two-dimensional (2D) materials. For systems that need
large simulation cells --- point defects, heterostructures, moir\'e superlattices with thousands
of atoms --- the Brillouin zone of the supercell is small, and it is natural to evaluate the
polarizability, the dielectric function and the screened interaction at a single wavevector,
the $\Gamma$ point of the supercell. The real-space resolution-of-the-identity (RI-RS)
formulation of $GW$ in \cpk{} is designed for this regime: the Green's function and the
polarizability are element-wise products on an atom-centred real-space grid, every contraction
with the auxiliary basis goes through sphere-local fit coefficients, and the cost grows with the
number of atoms rather than with the number of $k$-points.

The large-cell code in \cpk{} reconstructs $\chi(\bk)$, $\varepsilon(\bk)$ and $W(\bk)$ from
$\Gamma$ on a $k$-mesh \texttt{KPOINTS\_W}~$n_1\,n_2\,n_3$ of the supercell's own small
Brillouin zone and takes the bare Coulomb matrix $V(\bk)$ from a finite real-space image sum.
An even mesh is half-shifted and never samples $\bq=0$; an odd mesh does. The obvious $\Gamma$-only
setting, \texttt{KPOINTS\_W 1 1 1}, therefore samples exactly the wavevector at which the image
sum of the bare Coulomb operator between two charged auxiliary Gaussians diverges with the
summation range. It does not merely sample badly: it produces a large spurious additive
contribution to $W$. On a 392-atom hexagonal boron nitride (h-BN) slab the correlation
self-energy $\Sigma^c$ of the valence band maximum moved from $+2.195$ to $-24.043$~eV and that
of the conduction band minimum from $-3.283$ to $+20.237$~eV (opposite signs, because
$\Sigma^c=-G\,W$ and the two branches of $G$ carry opposite sign), while the exchange
self-energy $\Sigma^x$ and $v^{xc}$ stayed bit-identical.

This work makes \texttt{KPOINTS\_W 1 1 1} well defined. The divergence is shown to sit in one
direction of the auxiliary basis, the charge channel, which is removed from the sampled $V(0)$
and put back analytically through one additional basis function, the normalized constant. Its
contribution to the correlation part of the screened interaction (the ``head'' and the
``wings'') is known in closed form as a function of $\bq$ and is integrated over the whole
small Brillouin zone of the supercell, following the head-and-wing construction of
J.~Wilhelm for full-$k$ RI-RS $GW$ (Ref.~\onlinecite{wilhelm_fullk}, Sec.~VI; \cpk{}
PR~\#6447~\cite{pr6447}). Because the regular path is $\Gamma$ only and its sampled head is
removed exactly, nothing has to be subtracted.

Section~\ref{sec:theory2d} gives every equation of the scheme for a 2D slab in the order of
evaluation; Sec.~\ref{sec:theory3d} gives the changes for 3D bulk cells, where the Coulomb
matrix is an Ewald sum and the correction is applied in the orbital basis.
Section~\ref{sec:impl} describes the implementation in \cpk{}, Sec.~\ref{sec:results} the
results, and Sec.~\ref{sec:discussion} discusses why the old $k$-mesh large-cell scheme is not
a reference for 2D systems and why Gaussian form factors of the auxiliary charges can be
neglected.

%======================================================================================
\section{Theory for two-dimensional slabs}
\label{sec:theory2d}
%======================================================================================

\subsection{Setting and notation}
\label{sec:setting}

The simulation cell is a layer with in-plane lattice vectors $\mathbf a_1,\mathbf a_2$,
periodic in $x$ and $y$ only, with a non-periodic direction $z$. Its layer area is
\begin{equation}
  A = |\mathbf a_1\times\mathbf a_2| ,
\end{equation}
and the in-plane reciprocal vectors $\mathbf b_1,\mathbf b_2$ satisfy
$\mathbf b_i\cdot\mathbf a_j = 2\pi\delta_{ij}$. The vacuum height never enters any equation
below: all Fourier transforms average over a layer cell, not over a three-dimensional cell that
mostly contains vacuum.

Indices: $\mu,\nu,\lambda,\sigma$ label atomic orbitals (AOs) $\varphi_\mu$; $P,Q,S$ label
auxiliary (RI) Gaussians $\phi_P$; $\ell,\ell'$ label points $\br_\ell$ of a real-space grid;
$i$ labels occupied and $a$ empty Kohn--Sham states; $n,m$ label any state; $\bR$ labels lattice
vectors; $\tau$ is imaginary time and $\omega$ imaginary frequency. $\alpha,\beta\in\{x,y,z\}$
are Cartesian directions and $\qhat$ an in-plane unit vector. All quantities are real: there is
one real $k$-point, $\Gamma$, and the single complex factor of the construction (the wing of $W$)
is split off analytically.

The RI metric is a short-ranged kernel $v_\omega$, the truncated Coulomb operator
$v_\tc(r) = \theta(r_c-r)/r$ in the calculations reported here. Its zero-wavevector value is
\begin{equation}
  \int\dd\br\, v_\omega(r) =
  \begin{cases}
    \pi/\omega^2 & \text{erfc-attenuated metric},\\
    2\pi r_c^2 & \text{truncated Coulomb metric},
  \end{cases}
  \label{eq:metric_head}
\end{equation}
and is referred to as the metric head.

Imaginary time and frequency are discretized on minimax grids $\{\tau_j\}$, $\{\omega_k\}$,
$j,k = 1,\dots,N$. The cosine and sine transforms are written with their minimax weights,
\begin{equation}
\begin{aligned}
  f(\ii\omega_k) &= \sum_j \gamma^{\cos}_{kj}\, f(\ii\tau_j),\qquad
  f(\ii\tau_j) = \sum_k \tilde\gamma^{\cos}_{jk}\, f(\ii\omega_k),\\
  f(\ii\omega_k) &= \sum_j \gamma^{\sin}_{kj}\, f(\ii\tau_j),
\end{aligned}
  \label{eq:transforms}
\end{equation}
where the first represents $f(\ii\omega) = 2\int_0^\infty\dd\tau\cos(\omega\tau)f(\ii\tau)$ for
an even function of $\tau$, the second its inverse, and the third
$f(\ii\omega) = 2\ii\int_0^\infty\dd\tau\sin(\omega\tau)f(\ii\tau)$ for an odd one.

\subsection{Atomic orbitals on the real-space grid}
\label{sec:grid}

Modified Lebedev atomic grids are placed at every atom and their union is the global grid
$\{\br_\ell\}_{\ell=1}^{N_\mathrm{grid}}$, all points in the home cell. Because there is only one
$k$-point, the AO values on the grid are the $\Gamma$ Bloch sums over periodic images,
\begin{equation}
  \Phi_{\mu}(\br_\ell) = \sum_{\bR} \varphi_\mu(\br_\ell - \mathbf A_\mu - \bR),
  \label{eq:phi_grid}
\end{equation}
where $\mathbf A_\mu$ is the center of $\varphi_\mu$. The matrix $\Phi_{\ell\mu}$ is sparse:
$\Phi_{\ell\mu}$ is kept only if $|\br_\ell-\mathbf A_\mu-\bR|$ is within the AO radius for some
image. At the $14\times14$ h-BN cell 24.9\,\% of its elements are nonzero.

\subsection{Real-space RI: the fit coefficients $Z_{\ell P}$}
\label{sec:rirs}

RI-RS expresses products of AOs through the grid. Define the grid metric and the grid--RI
right-hand side,
\begin{align}
  D_{\ell\ell'} &= \Big[\sum_\mu \Phi_\mu(\br_\ell)\Phi_\mu(\br_{\ell'})\Big]^2 ,
  \label{eq:D}\\
  d_{\ell P} &= \sum_{\mu\nu}\Phi_\mu(\br_\ell)\Phi_\nu(\br_\ell)\,(\mu\nu|P)_\omega ,
  \label{eq:d}
\end{align}
with
\begin{equation}
  (\mu\nu|P)_\omega = \sum_{\bR\bR'}\!\int\!\!\int\dd\br\,\dd\br'\,
  \varphi_{\mu}^{\bR}(\br)\varphi_{\nu}^{\bR'}(\br)\, v_\omega(|\br-\br'|)\,\phi_P(\br') .
  \label{eq:3c}
\end{equation}
Writing $B_{(\mu\nu),\ell} = \Phi_\mu(\br_\ell)\Phi_\nu(\br_\ell)$ one has
$D = B^\Tr B$ and $d_{\cdot P} = B^\Tr g_P$ with $g_{(\mu\nu),P} = (\mu\nu|P)_\omega$, so the
normal equations
\begin{equation}
  \sum_{\ell'} D_{\ell\ell'} Z_{\ell' P} = d_{\ell P}
  \label{eq:fit}
\end{equation}
are the least-squares solution of
$\sum_\ell \Phi_\mu(\br_\ell)\Phi_\nu(\br_\ell) Z_{\ell P}\approx (\mu\nu|P)_\omega$: the
coefficient $Z_{\ell P}$ is a quadrature weight times the value at $\br_\ell$ of the metric
potential of $\phi_P$. The solve is balanced and regularized,
\begin{equation}
\begin{gathered}
  j_\ell = \frac{1}{\max(D_{\ell\ell},\sqrt{\epsilon_\mathrm{floor}})},\\
  \big[j_\ell D_{\ell\ell'} j_{\ell'} + \lambda\,\delta_{\ell\ell'}\big] Z'_{\ell'P} = j_\ell d_{\ell P},
  \qquad Z_{\ell P} = j_\ell Z'_{\ell P},
\end{gathered}
  \label{eq:fit_reg}
\end{equation}
with the Tikhonov parameter $\lambda = 10^{-8}$ and $\epsilon_\mathrm{floor}=10^{-16}$.
Because $d_{\ell P}$ decays away from the center of $\phi_P$, Eq.~\eqref{eq:fit} is solved
independently for each atom $P$ on the grid points inside a sphere around it (radius
$r_\mathrm{RI}+r_\mathrm{AO}$), by a dense Cholesky factorization of the sphere's block of
$D$. This sphere locality is what makes RI-RS linear scaling; it is valid precisely because
every column $d_{\ell P}$ is localized, a point that becomes important in Sec.~\ref{sec:Zl0}.

Every subsequent contraction between the grid and the RI basis uses the same matrix:
\begin{equation}
  X_{PQ} = \sum_{\ell\ell'} Z_{\ell P}\, X_{\ell\ell'}\, Z_{\ell' Q},\qquad
  X_{\ell\ell'} = \sum_{PQ} Z_{\ell P}\, X_{PQ}\, Z_{\ell' Q}.
  \label{eq:ZXZ}
\end{equation}

\subsection{Green's functions and polarizability at $\Gamma$}
\label{sec:chi}

With Kohn--Sham coefficients $C_{\mu n}$, eigenvalues $\varepsilon_n$ and the Fermi level
$\varepsilon_F$ placed mid-gap, the occupied and empty Green's functions in imaginary time are
\begin{align}
  G^\occ_{\mu\nu}(\ii|\tau|) &= \sum_{i}^{\occ} C_{\mu i}\,\ee^{-|(\varepsilon_i-\varepsilon_F)\tau|}\,C_{\nu i},\nonumber\\
  G^\vir_{\mu\nu}(\ii|\tau|) &= \sum_{a}^{\vir} C_{\mu a}\,\ee^{-|(\varepsilon_a-\varepsilon_F)\tau|}\,C_{\nu a},
  \label{eq:G_ao}\\
  G^{\occ/\vir}_{\ell\ell'}(\ii|\tau|) &= \sum_{\mu\nu}\Phi_\mu(\br_\ell)\, G^{\occ/\vir}_{\mu\nu}(\ii|\tau|)\,\Phi_\nu(\br_{\ell'}),
  \label{eq:G_grid}
\end{align}
each built as $(\tilde C)(\tilde C)^\Tr$ with $\tilde C_{\mu n} = C_{\mu n}\ee^{-|(\varepsilon_n-\varepsilon_F)\tau|/2}$.
The independent-particle polarizability is an element-wise (Hadamard) product on the grid,
contracted to the RI basis with Eq.~\eqref{eq:ZXZ},
\begin{equation}
\begin{aligned}
  \chi_{\ell\ell'}(\ii\tau) &= G^\occ_{\ell\ell'}(\ii|\tau|)\,G^\vir_{\ell\ell'}(\ii|\tau|),\\
  \chi_{PQ}(\ii\tau) &= \sum_{\ell\ell'} Z_{\ell P}\,\chi_{\ell\ell'}(\ii\tau)\,Z_{\ell'Q}.
\end{aligned}
  \label{eq:chi}
\end{equation}
The grid product is evaluated in panels of grid rows; a geometric cutoff on $|\br_\ell-\br_{\ell'}|$
limits the pairs. At $\Gamma$, $\chi_{\ell\ell'}$ is a Gram matrix of occupied--empty products
and therefore exactly positive semidefinite; no clipping of negative eigenvalues is needed. It is
transformed to imaginary frequency with the cosine weights of Eq.~\eqref{eq:transforms},
$\chi_{PQ}(\ii\omega_k) = \sum_j\gamma^{\cos}_{kj}\chi_{PQ}(\ii\tau_j)$.

\subsection{The screened interaction at $\Gamma$: the problem and the plan}
\label{sec:plan}

$W$ is needed at one point only, $\bq=0$. There the bare Coulomb matrix between charged
auxiliary functions diverges. The scheme handles this in three steps:
\begin{enumerate}
\item The divergence sits in one direction of the auxiliary basis, the charge vector $c$. It is
  removed from $V(0)$ as a single constant, $V(0)\to V(0)-D_\mathrm{box}\,cc^\Tr$
  (Sec.~\ref{sec:V}).
\item $\varepsilon$ and $W$ are computed from this finite $V(0)$ in the usual way
  (Sec.~\ref{sec:W}). What is now missing from $W$ is exactly the screened interaction in the
  charge channel.
\item That missing piece is added back analytically: the charge channel is represented by the
  constant function $\phi_0 = 1/\sqrt A$, its contribution to $W$ (the ``head'' and ``wings'')
  is known in closed form as a function of $\bq$, and it is integrated over the whole small
  Brillouin zone of the supercell (Sec.~\ref{sec:charge}). Nothing has to be subtracted,
  because step~1 already removed it from the sampled point.
\end{enumerate}

\subsection{Step 1: the bare Coulomb matrix at $\bq=0$}
\label{sec:V}

\subsubsection{What is computed}

The Coulomb matrix between two auxiliary functions $\phi_P$ and $\phi_Q$ in a periodic layer is
the interaction of $\phi_P$ in the home cell with $\phi_Q$ in \emph{every} cell. The code sums
over a finite box of cells,
\begin{equation}
  V^\lat_{PQ}(0) = \sum_{\bR\in\mathcal B}\underbrace{\int\!\!\int\dd\br\,\dd\br'\,
  \frac{\phi_P(\br)\,\phi_Q(\br'-\bR)}{|\br-\br'|}}_{(P0|Q\bR)},
  \label{eq:Vlat}
\end{equation}
\begin{equation}
  \mathcal B = \{n_1\mathbf a_1+n_2\mathbf a_2 : |n_1|,|n_2|\le m\},
  \label{eq:box}
\end{equation}
with $m = (3^{s-1}-1)/2$ and $s$ the input \texttt{SIZE\_LATTICE\_SUM} ($s=4$: $m=13$, i.e.\
$27\times27$ cells). There are no images in the direction normal to the layer.

\subsubsection{Far images see only the charge}

Take one term $(P0|Q\bR)$ with $|\bR|$ much larger than the extent of the two Gaussians.
Seen from far away, $\phi_Q$ looks like a point charge. Its charge is its integral,
\begin{equation}
  c_Q = \int\dd\br\,\phi_Q(\br),
  \label{eq:charge}
\end{equation}
and the same for $\phi_P$. The multipole expansion of the term is therefore
\begin{equation}
  (P0|Q\bR) = \frac{c_P\,c_Q}{|\bR|} + \frac{(\text{dipole terms})}{|\bR|^2} + \frac{(\ldots)}{|\bR|^3} + \cdots
  \label{eq:multipole}
\end{equation}
Summed over the box, the dipole terms cancel between $\bR$ and $-\bR$ and the $|\bR|^{-3}$ and
higher terms converge. The only part of the sum that does not settle down is the first one:
\begin{equation}
  V^\lat_{PQ}(0) = c_P\,c_Q\,\sum_{\bR\in\mathcal B}\frac{1}{|\bR|}
  \;+\;\big(\text{converges as the box grows}\big).
  \label{eq:split}
\end{equation}

\subsubsection{Why the first part diverges}

In two dimensions the number of lattice points in a ring of radius $R$ and width $\dd R$ is
$2\pi R\,\dd R/A$ (one point per area $A$). Each of them contributes $1/R$, so the ring
contributes $2\pi\,\dd R/A$ --- the same for every ring. The sum therefore grows
\emph{linearly} with the box radius and never converges:
\begin{equation}
  \sum_{\bR\in\mathcal B}\frac{1}{|\bR|}\;\approx\;\frac{1}{A}\int_{\mathcal P}\frac{\dd^2 r}{|\br|}
  \;\propto\;\frac{2\pi}{A}\,R_\mathrm{box}.
\end{equation}
On the h-BN cell this term is $7.4\times10^{4}$~Ha in the charge direction (defined next),
against a few Ha for everything else.

\subsubsection{The charge channel}

Collect the charges into one vector $c = (c_1,\dots,c_{N_\mathrm{RI}})$ and normalize it,
$\hat c = c/\|c\|$. Equation~\eqref{eq:split} says the divergent part of the matrix is
$c\,c^\Tr$ times a number: it is a \emph{rank-one} matrix, nonzero only along the single
direction $\hat c$ of the auxiliary basis. We call that direction the \emph{charge channel}.
Any matrix can be split into its charge-channel part and the rest,
\begin{equation}
  V = \underbrace{(\hat c^\Tr V\hat c)\,\hat c\hat c^\Tr}_{\text{charge channel}}
  + \underbrace{\big(V - (\hat c^\Tr V\hat c)\,\hat c\hat c^\Tr\big)}_{\text{rest}},
\end{equation}
and the whole problem of $\bq=0$ is that the number $\hat c^\Tr V\hat c$ is infinite. Two
remarks matter later. First, not every auxiliary function is charged: an $s$-type
Gaussian is, a $p$-type one is not, and a combination like $\phi_1 - (c_1/c_2)\phi_2$ is
neutral. Second, the charge channel also contains a \emph{finite} part --- the interaction of a
charged function with its own neighbourhood, the second term of Eq.~\eqref{eq:split} projected
onto $\hat c$ --- which is physical and must be kept. Only the divergent constant may go.

\subsubsection{$D_\mathrm{box}$: the constant that is removed}

The divergent number is the lattice sum of $1/|\bR|$ over the box. Replacing the sum by the
integral over the region $\mathcal P$ that the summed cells cover --- the parallelogram
$\{u_1\mathbf a_1+u_2\mathbf a_2 : |u_1|,|u_2|\le F\}$ with $F = 3^{s-1}/2$, one cell per
lattice point --- defines
\begin{equation}
  D_\mathrm{box} = \frac{1}{A}\int_{\mathcal P}\frac{\dd^2 r}{|\br|} .
  \label{eq:Dbox}
\end{equation}
In polar coordinates $\dd^2r = r\,\dd r\,\dd\theta$, the $1/r$ cancels, and the radial
integral is just the distance $r_{\max}(\theta)$ from the origin to the edge of the
parallelogram in direction $\theta$:
\begin{equation}
  D_\mathrm{box} = \frac{1}{A}\int_0^{2\pi}\dd\theta\; r_{\max}(\theta)
  = \frac{1}{A}\oint_{\partial\mathcal P}\frac{\br\times\dd\br}{|\br|},
  \label{eq:Dbox_edge}
\end{equation}
a one-dimensional integral along the four edges of the box, done by Gauss--Legendre quadrature.
It is a pure number fixed by the cell shape and the box size: $2.80$ for the $6$~\AA\ water cell
at $s=3$, $1.52$ for the h-BN cell at $s=4$.

\subsubsection{The Coulomb matrix used for $W$}

Subtracting the constant along the charge direction, with the \emph{unnormalized} charges so
that the prefactor of Eq.~\eqref{eq:split} is matched,
\begin{equation}
  V_{PQ}(0) = V^\lat_{PQ}(0) - D_\mathrm{box}\, c_P\, c_Q .
  \label{eq:Vreg}
\end{equation}
This removes exactly the first term of Eq.~\eqref{eq:split} and nothing else. What is left in
the charge channel is small and finite: $\hat c^\Tr V\hat c$ goes from $796$ to $8.9$~Ha for the
water cell and from $7.4\times10^{4}$ to $-139$~Ha for h-BN. In reciprocal space this is the
familiar ``omit the $\bG=0$ term'' of the Coulomb sum; Eq.~\eqref{eq:Vreg} is its real-space
form.

What one must \emph{not} do is project the charge channel out altogether,
$V\to\Pi V\Pi$ with $\Pi = 1-\hat c\hat c^\Tr$. That would also discard the finite part of the
charge channel (the second remark above), which step~3 does not put back, and it was measured
to inflate every $\Sigma^c$ matrix element by a factor of 5--9 on the water test.

\subsection{Step 2: $\varepsilon$ and $W$ from the finite $V(0)$}
\label{sec:step2}
\label{sec:W}

With the RI metric $M_{PQ} = (P|v_\omega|Q)$ and $V$ from Eq.~\eqref{eq:Vreg},
\begin{align}
  \varepsilon(\ii\omega) &= 1 + V^{1/2} M^{-1}\chi(\ii\omega) M^{-1} V^{1/2},
  \label{eq:eps}\\
  W^c(\ii\omega) &= V^{1/2}\big(\varepsilon^{-1}(\ii\omega) - 1\big)V^{1/2},
  \label{eq:Wc}\\
  W^{\MIC}(\ii\omega) &= M^{-1}\, W^c(\ii\omega)\, M^{-1}.
  \label{eq:Wmic}
\end{align}
This is the only $k$-point, so $W^\MIC$ is Eq.~\eqref{eq:Wmic} with weight one. (The code's
two-mesh extrapolation must be switched off for this; with it, $\Gamma$ would enter $W$ with
weight $1/(1-\sqrt2)$ in some matrix elements and weight 0 in others.) $W^\MIC$ is transformed
to imaginary time and brought to the grid,
\begin{equation}
  W^{\MIC}_{\ell\ell'}(\ii\tau) = \sum_{PQ} Z_{\ell P}\, W^{\MIC}_{PQ}(\ii\tau)\, Z_{\ell'Q}.
  \label{eq:W_grid}
\end{equation}

\subsection{Step 3: the charge channel, added back analytically}
\label{sec:charge}

\subsubsection{The constant function and its auxiliary charges}
\label{sec:Zl0}

The charge channel is represented by one extra basis function, the normalized constant,
\begin{equation}
  \phi_0(\br) = \frac{1}{\sqrt A},\qquad V_{00}(\bq) = \frac{2\pi}{q}.
  \label{eq:phi0}
\end{equation}
$2\pi/q$ is the exact two-dimensional Coulomb head; it is finite for every $q>0$ and the
$1/q$ singularity is integrable in two dimensions, which is why integrating over the Brillouin
zone works where sampling at $q=0$ does not.

At a small but finite $\bq$ the charge channel is the Bloch-modulated constant
$\ee^{\ii\bq\cdot\br}/\sqrt A$, and a fitted density $\rho = \sum_P \rho_P\phi_P$ couples to
it through its Fourier component,
\begin{equation}
\begin{aligned}
  \langle\phi_{0,\bq}|\rho\rangle &= \frac{\hat\rho(\bq)}{\sqrt A}
  = \frac{1}{\sqrt A}\sum_P \rho_P\,\hat\phi_P(\bq)\\
  &\simeq \frac{1}{\sqrt A}\sum_P \rho_P\, c_P\,\ee^{-\ii\bq\cdot\bR_P},\qquad
  c_P = \int\dd\br\,\phi_P(\br),
\end{aligned}
  \label{eq:Zl0}
\end{equation}
with $\bR_P$ the centre of the auxiliary Gaussian $\phi_P$. The last step is the point-charge
limit $q_\mathrm{max}\sigma_P\ll1$ of the Gaussian form factor; only $s$-type functions carry a
charge. So the charge channel needs no fit of its own: the auxiliary charges $c_P$ and the
atomic positions are everything, and the correction is applied in the auxiliary basis
(Sec.~\ref{sec:DeltaW}), reaching the grid through the same $Z$ as every other part of $W$.

An earlier version gave $\phi_0$ grid coefficients $Z_{\ell0}$ from the fit of
Eq.~\eqref{eq:fit} with the constant as the target (solved globally by conjugate gradients, since
its right-hand side does not decay) and applied the correction per grid pair,
$Z_{\ell0}Z_{\ell'0}K_\head(\bR_{\ell\ell'})$. That is numerically unusable. The RI-RS
interpolation functions $\xi_\ell$ far from the atoms carry fit weights of $10^4$--$10^8$ of
alternating sign (on the 3$\times$3 h-BN test cell, $\max|Z_{\ell0}| = 2\times10^4$ at points
with a density of $10^{-9}$). They cancel exactly in linear functionals of a density,
$\sum_\ell\rho(\br_\ell)Z_{\ell0}$ gives $1/\sqrt A$ for a diagonal pair and $3\times10^{-6}$
for a transition density, but not once each point is weighted by a different kernel value
$K(\bR_{\ell\ell'})$: the odd wing came out wrong by orders of magnitude and split
symmetry-equivalent states by more than 1~eV, the even head survived only on large cells where
the kernel is nearly flat over neighbouring points.

\subsubsection{The neutral Gaussians}
\label{sec:neutral}

The Gaussians are made neutral by projecting out the charge direction,
\begin{equation}
  \Pi = 1 - \hat c\hat c^\Tr,\qquad V_b = \Pi\, V^\lat(0)\,\Pi ,
  \label{eq:Pi}
\end{equation}
so that $V_b$ is finite (neutral functions interact as $R^{-3}$). The two maps between a fitted
grid density and the neutral basis are
\begin{equation}
  F^\dens = M^{-1}\,\Pi\, V_b^{1/2},\qquad
  F^\ext = M^{-1}\, V^\lat(0)\,\Pi\, V_b^{-1/2},
  \label{eq:F}
\end{equation}
the first for the density side of $\varepsilon$, the second for the potential side of $W$. We
also write $G^\ext = M F^\ext = V^\lat(0)\,\Pi\,V_b^{-1/2}$ for the potential-side map in the
representation of $W^c$.

\subsubsection{The response in the charge channel}
\label{sec:response}

The constant function couples to the electrons through the transition dipoles. Under periodic
boundary conditions the position operator is ill defined, so they are taken in the velocity
gauge,
\begin{equation}
  \mu^\alpha_{ai} = \frac{\langle a|[H,r^\alpha]|i\rangle}{\varepsilon_a-\varepsilon_i},\qquad
  [H,r^\alpha] = -\partial_\alpha + [V_\mathrm{nl},r^\alpha],
  \label{eq:dipole}
\end{equation}
and checked against the Thomas--Reiche--Kuhn sum rule
\begin{equation}
  \frac{2}{n_\mathrm{spin}}\sum_{\sigma}\sum_{ia}2\,(\varepsilon_a-\varepsilon_i)\,|\mu^\alpha_{ai}|^2 = N_\mathrm{el},
  \label{eq:TRK}
\end{equation}
which a finite basis can only undershoot (h-BN: $0.64/0.64/0.41$). The head and the wings of
the polarizability in imaginary time are
\begin{align}
  Y_{\alpha\beta}(\ii\tau) &= \frac{2}{n_\mathrm{spin}A}\sum_\sigma\sum_{ia}\mu^\alpha_{ai}\mu^\beta_{ai}\,\ee^{-(\varepsilon_a-\varepsilon_i)\tau},
  \label{eq:Y}\\
  X_{\ell\alpha}(\ii\tau) &= \frac{2}{n_\mathrm{spin}\sqrt A}\sum_\sigma\sum_{ia}\psi_i(\br_\ell)\psi_a(\br_\ell)\,\mu^\alpha_{ai}\,\ee^{-(\varepsilon_a-\varepsilon_i)\tau},
  \label{eq:X}
\end{align}
and are cosine-transformed to imaginary frequency like $\chi$. $Y(\ii\omega\to0)$ is the
polarizability per area of the layer (h-BN: 1.35~bohr).

\subsubsection{Inverting $\varepsilon$ in the extended basis}
\label{sec:blockinv}

With the constant as index 0, the dielectric matrix at small $\bq$ is a $2\times2$ block matrix,
\begin{equation}
  \varepsilon(\bq,\ii\omega) = \begin{pmatrix} h & w^\dagger\\ w & B\end{pmatrix},
  \label{eq:blocks}
\end{equation}
\begin{equation}
  h = 1 + 2\pi q\,\qhat^\Tr Y\qhat,\qquad
  w_P = -\ii\sqrt{2\pi q}\,(E\qhat)_P,
  \label{eq:hw}
\end{equation}
\begin{equation}
  B(\ii\omega) = 1 + F^{\dens\,\Tr}\chi(\ii\omega)F^\dens,\quad
  E(\ii\omega) = F^{\dens\,\Tr}Z^\Tr X(\ii\omega).
  \label{eq:E}
\end{equation}
$B$ is the ordinary Gaussian dielectric matrix at $\Gamma$ in the neutral basis; $h$ and $w$
are the new head and wings. $B$ needs no new grid contraction: $Z F^\dens$ maps the grid to the
neutral basis, so $F^{\dens\,\Tr}(Z^\Tr\chi_{\ell\ell'}Z)F^\dens$ contains exactly the
$\chi_{PQ}(\Gamma)$ of Eq.~\eqref{eq:chi}. Inverting the block matrix and keeping the head and
the wings of $W^c = V^{1/2}(\varepsilon^{-1}-1)V^{1/2}$ gives closed forms. All of the $\bq$
dependence collapses into one scalar,
\begin{equation}
  a(\qhat,\ii\omega) = \qhat^\Tr S\qhat,\qquad
  S = Y - E^\Tr B^{-1}E,
  \label{eq:a}
\end{equation}
\begin{equation}
  T_{P\alpha} = \big[V^\lat(0)\,\Pi\,V_b^{-1/2}\,B^{-1}E\big]_{P\alpha},
  \label{eq:ST}
\end{equation}
($S$ is the $3\times3$ head with the body screening folded in, $T$ the wing vector in the
auxiliary basis --- $M F^\ext B^{-1}E = G^\ext B^{-1}E$, i.e.\ in the same representation as
$W^c$ of Eq.~\eqref{eq:Wc} --- both obtained by linear solves with $B$), and then
\begin{align}
  W^{c}_{\head}(\bq,\ii\omega) &= -\frac{4\pi^2 a}{1+2\pi q a},
  \label{eq:headwing}\\
  W^{c}_{\wing,P}(\bq,\ii\omega) &= \frac{\ii\,2\pi\sum_\alpha\hat q_\alpha T_{P\alpha}}{1+2\pi q a}.
  \label{eq:wing}
\end{align}
Both are finite at $q\to0$: the head of $\varepsilon^{-1}$ is proportional to $q$ and cancels
the $1/q$ of $V_{00}$, the wing is proportional to $\sqrt q$ and cancels $1/\sqrt q$. The head is
a number, the wing is a vector over the auxiliary functions and linear in $\qhat$, which is what
allows the angular integral to be done once into a table. The explicit factors of $\ii$ cancel
in the Schur complement $h - w^\dagger B^{-1}w = 1+2\pi q a$ and reappear in $W$ only as a
quantity plus its own conjugate, so every stored array is real.

\subsubsection{Integral over the Brillouin zone and the correction to $W$}
\label{sec:DeltaW}

The head and wings are now integrated over the whole small Brillouin zone of the supercell,
with the phase $\ee^{-\ii\bq\cdot\bR}$ of Eq.~\eqref{eq:Zl0} that turns $W(\bq)$ into a
function of the separation $\bR = \MIC(\bR_P-\bR_Q)$ of the atoms carrying two auxiliary
functions. Because $a$ is even and the wing odd in
$\qhat$, only $\cos$ survives in the head and only $\sin$ in the wing, and both results are
real:
\begin{align}
  K_\head(\bR,\ii\omega) &= \frac{A}{(2\pi)^2}\int_{\mathrm{BZ}}\dd^2 q\;
  \Big[-\frac{4\pi^2 a}{1+2\pi q a}\Big]\cos(\bq\cdot\bR),
  \label{eq:Kh}\\
  K_{\wing,\alpha}(\bR,\ii\omega) &= \frac{A}{(2\pi)^2}\int_{\mathrm{BZ}}\dd^2 q\;
  \Big[\frac{2\pi\,\hat q_\alpha}{1+2\pi q a}\Big]\sin(\bq\cdot\bR).
  \label{eq:Kw}
\end{align}
The prefactor $A/(2\pi)^2$ makes this the Brillouin-zone \emph{average}, the continuum limit of
$\frac1{N_k}\sum_k$. The integrals are done in polar coordinates (24 Gauss--Legendre directions
per zone edge, 32 radial points), tabulated on a $128\times128$ mesh of $\bR$ over the
minimum-image cell and interpolated bilinearly at the atom-pair separations. The correction
to the screened interaction in the auxiliary basis is
\begin{equation}
\begin{aligned}
  \Delta W_{PQ}(\ii\omega) ={}& \frac{c_P c_Q}{A}\,K_\head(\bR_{PQ},\ii\omega)\\
  &+ \frac{1}{\sqrt A}\sum_\alpha\big[T_{P\alpha}c_Q - c_P T_{Q\alpha}\big]
  K_{\wing,\alpha}(\bR_{PQ},\ii\omega),
\end{aligned}
  \label{eq:DeltaW}
\end{equation}
symmetric in $P\leftrightarrow Q$ ($K_\head$ is even and $K_\wing$ odd in $\bR$, the wing
bracket antisymmetric). It is transformed to imaginary time with the cosine weights of the
minimax grid and added to $W^c_{PQ}(\ii\tau)$ \emph{before} Eq.~\eqref{eq:Wmic}, so that it
reaches the grid pairs through $M^{-1}$ and $Z$ exactly like the regular $W$. Its head term
is a uniform potential $K_\head(0)/A$ in real space ($-0.29$~eV for h-BN). The cost is
negligible: one table lookup per auxiliary pair and time point.

\subsection{Exchange self-energy}
\label{sec:Sx}

The exchange part uses the truncated Coulomb operator $v_\tc$ with a cutoff radius equal to
half the minimum cell dimension (15.2~\AA\ for the h-BN cell), which is finite without any
treatment of the charge channel and is unaffected by everything above:
\begin{align}
  P_{\mu\nu} &= \sum_i^{\occ} C_{\mu i}C_{\nu i},\qquad
  P_{\ell\ell'} = \sum_{\mu\nu}\Phi_\mu(\br_\ell)P_{\mu\nu}\Phi_\nu(\br_{\ell'}),
  \label{eq:Sx_parts}\\
  V^\tc_{\ell\ell'} &= \sum_{PQ}Z_{\ell P}\big[M^{-1}V^\tc M^{-1}\big]_{PQ}Z_{\ell'Q},
  \label{eq:Vtc_grid}\\
  \Sigma^x_{\lambda\sigma} &= -\sum_{\ell\ell'}\Phi_\lambda(\br_\ell)\,P_{\ell\ell'}\,V^\tc_{\ell\ell'}\,\Phi_\sigma(\br_{\ell'}).
  \label{eq:Sx}
\end{align}

\subsection{Correlation self-energy and quasiparticle energies}
\label{sec:Sc}

The correlation self-energy is an element-wise product on the grid of the Green's function with
the \emph{full} screened interaction, the regular part of Eq.~\eqref{eq:W_grid} plus the
charge-channel correction of Eq.~\eqref{eq:DeltaW}, in the two time branches:
\begin{align}
  \Sigma^{c}_{\ell\ell'}(-\ii|\tau|) &= -\,G^\occ_{\ell\ell'}(\ii|\tau|)\,
  \Big[W^\MIC_{\ell\ell'}(\ii\tau) + \Delta W_{\ell\ell'}(\ii\tau)\Big],
  \label{eq:Sc_neg}\\
  \Sigma^{c}_{\ell\ell'}(+\ii|\tau|) &= +\,G^\vir_{\ell\ell'}(\ii|\tau|)\,
  \Big[W^\MIC_{\ell\ell'}(\ii\tau) + \Delta W_{\ell\ell'}(\ii\tau)\Big],
  \label{eq:Sc_pos}\\
  \Sigma^{c}_{\lambda\sigma}(\pm\ii|\tau|) &= \sum_{\ell\ell'}\Phi_\lambda(\br_\ell)\,
  \Sigma^{c}_{\ell\ell'}(\pm\ii|\tau|)\,\Phi_\sigma(\br_{\ell'}).
  \label{eq:Sc_ao}
\end{align}
In the implementation $W^\MIC_{\ell\ell'}$ is built panel by panel as $Z\,W^\MIC_{PQ}Z^\Tr$,
with $\Delta W$ already inside $W^\MIC_{PQ}$ (Eq.~\eqref{eq:DeltaW}), and the two Hadamard
products with $G^\occ$ and $G^\vir$ share that panel. The panels are restricted to the pairs
within the geometric cutoff of the $\Sigma^c$ panels, which is exact because the Hadamard
product with the short-ranged $G$ vanishes outside it.

The self-energy is projected onto the Kohn--Sham states,
$\Sigma^c_{nn}(\pm\ii|\tau|) = \sum_{\lambda\sigma}C_{\lambda n}\Sigma^c_{\lambda\sigma}(\pm\ii|\tau|)C_{\sigma n}$,
split into its even and odd parts in $\tau$ and transformed to imaginary frequency with the
cosine and sine weights of Eq.~\eqref{eq:transforms},
\begin{equation}
\begin{aligned}
  \Sigma^c_{nn}(\ii\omega_k) ={}& \sum_j\gamma^{\cos}_{kj}\,
  \frac{\Sigma^c_{nn}(\ii\tau_j)+\Sigma^c_{nn}(-\ii\tau_j)}{2}\\
  &+ \sum_j\gamma^{\sin}_{kj}\,
  \frac{\Sigma^c_{nn}(\ii\tau_j)-\Sigma^c_{nn}(-\ii\tau_j)}{2}.
\end{aligned}
  \label{eq:Sc_freq}
\end{equation}
$\Sigma^c_{nn}(\ii\omega)$ is continued to the real axis by a Pad\'e approximant fitted to
the imaginary-frequency points, and the quasiparticle energies follow from
\begin{equation}
  \varepsilon^{GW}_n = \varepsilon_n + \mathrm{Re}\,\Sigma^c_{nn}(\varepsilon^{GW}_n)
  + \Sigma^x_{nn} - v^{xc}_{nn},
  \label{eq:QP}
\end{equation}
solved by a Newton search for the crossing of the quasiparticle equation. For $G_0W_0$ this is
done once with $G$ built from the Kohn--Sham eigenvalues; for ev$GW_0$ the Green's function is
rebuilt from the quasiparticle energies and Eqs.~\eqref{eq:G_ao}--\eqref{eq:Sc_freq} are
iterated to self-consistency with $W$ frozen.

%======================================================================================
\section{Theory for three-dimensional bulk cells}
\label{sec:theory3d}
\label{sec:bulk}
%======================================================================================

\subsection{Substitutions}
\label{sec:subst3d}

The same construction applies to a bulk cell (\texttt{PERIODIC XYZ}) with three substitutions.
The measure is the cell volume $\Omega = |\det(\mathbf a_1,\mathbf a_2,\mathbf a_3)|$,
$A\to\Omega$ in Eqs.~\eqref{eq:Zl0}, \eqref{eq:Y}, \eqref{eq:X} and \eqref{eq:DeltaW}, the
reciprocal vectors are the three $\mathbf b_i$ with $\mathbf b_i\cdot\mathbf a_j=2\pi\delta_{ij}$,
and the charge mode is $\phi_0 = \ee^{\ii\bq\cdot\br}/\sqrt\Omega$ with $v_{00} = 4\pi/q^2$.
Everything else --- the dipoles, the response $Y$, $X$, $B$, $E$, the block inverse and the
transform $\ii\omega\to\ii\tau$ --- is shared with the 2D case.

\subsection{The Coulomb matrix: Ewald sum without $\bG=0$}
\label{sec:ewald}

The finite box sum of Eq.~\eqref{eq:Vlat} cannot be used in three dimensions, not even after
the subtraction of Eq.~\eqref{eq:Vreg}: between two \emph{neutral} (dipolar) functions the sum
of $1/r$ over a 3D box is conditionally convergent, with a surface term
$\tfrac12 d_\alpha d_\beta\int_\mathrm{box}\partial_\alpha\partial_\beta r^{-1}$ that depends
on the shape of the box and does not vanish as the box grows (the Lorentz depolarization of a
finite sample). The matrix that the head-and-wing model assumes is the tin-foil limit, the
periodic Coulomb matrix with the $\bG=0$ Fourier coefficient of $1/r$ omitted, which is
obtained as an Ewald sum,
\begin{widetext}
\begin{equation}
  V_{PQ}(0) = \sum_{\mathbf{L}\in\mathrm{box}}\big(P\,0\,\big|\,\tfrac{\mathrm{erfc}(\omega r)}{r}\,\big|\,Q\,\mathbf{L}\big)
  + \frac{4\pi}{\Omega}\sum_{\bG\neq0}\frac{\ee^{-G^2/4\omega^2}}{G^2}\,\rho_P(\bG)\rho_Q(\bG)^*
  - \frac{\pi}{\Omega\,\omega^2}\,c_P c_Q,
  \qquad \rho_P(\bG) = \int\dd\br\,\phi_P(\br)\,\ee^{\ii\bG\cdot\br},
  \label{eq:Vewald}
\end{equation}
\end{widetext}
where $1/r = \mathrm{erfc}(\omega r)/r + \mathrm{erf}(\omega r)/r$, the erf part has the
Fourier transform $(4\pi/G^2)\,\ee^{-G^2/4\omega^2} = 4\pi/G^2 - \pi/\omega^2 + O(G^2)$, and
dropping its $4\pi/G^2$ singularity leaves $-\pi/\omega^2$ at $\bG=0$. (Here $\omega$ is the
Ewald splitting parameter, not an imaginary frequency.) The result is independent of $\omega$,
which checks both halves against each other; the code sets $\omega r_c = 8$ with
$r_c = \tfrac12(N-1)h_\mathrm{min}$ the coverage radius of the image box ($N = 3^{s-1}$ images
per direction, $h_\mathrm{min}$ the smallest cell height $2\pi/|\mathbf b_i|$), so that the erfc
tail is negligible for every pair the box contains ($\mathrm{erfc}(8)<10^{-28}$; the Gaussian
extents of the auxiliary functions use a few bohr of that margin). $\rho_P(\bG)$ is analytic
for Gaussians (one-dimensional moments
$m_l(p) = \int\dd x\,x^l\ee^{-\alpha x^2}\ee^{\ii p x}$ by the recursion
$m_l = (\ii p\,m_{l-1} + (l-1)m_{l-2})/2\alpha$). This $V(0)$ is finite and positive on the
charge direction, so no projector is needed: $\Pi = 1$, $V_b = V(0)$,
$F^\dens = M^{-1}V^{1/2}$ and $G^\ext = V^{1/2}$, so that Eq.~\eqref{eq:ST} becomes
\begin{equation}
  T = V(0)\,V^{-1/2}\,B^{-1}E = V^{1/2}B^{-1}E .
  \label{eq:T3d}
\end{equation}
The regular path of Sec.~\ref{sec:step2} uses the same matrix with nothing subtracted (the same
choice as the 3D branch of \cpk{}~\#6447~\cite{pr6447}).

\subsection{Head and wings}
\label{sec:hw3d}

In the extended basis the head is $h = 1 + 4\pi\,\qhat^\Tr Y\qhat$ and the wing
$w_P = -\ii\sqrt{4\pi}\,(E\qhat)_P$, both independent of $q$: the $q^2$ of the dipole response
cancels the $1/q^2$ of $v_{00}$ where in two dimensions $2\pi q\,a$ was left. The Schur
complement of the symmetrized dielectric matrix is
\begin{equation}
  \tilde\varepsilon(\qhat,\ii\omega) = 1 + 4\pi a(\qhat,\ii\omega),\qquad a = \qhat^\Tr S\qhat,
  \label{eq:epstilde}
\end{equation}
with $S$ from Eq.~\eqref{eq:a}, and the block inverse gives
\begin{align}
  W^c_{\head}(\bq,\ii\omega) &= \frac{4\pi}{q^2}\Big(\frac{1}{1+4\pi a}-1\Big)
  = -\frac{4\pi}{q^2}\,\frac{4\pi a}{1+4\pi a},
  \label{eq:headwing3d}\\
  W^c_{\wing,P}(\bq,\ii\omega) &= \frac{\ii\,(4\pi/q)\sum_\alpha\hat q_\alpha T_{P\alpha}}{1+4\pi a}.
  \label{eq:wing3d}
\end{align}
$1 + 4\pi a(\qhat,\ii\omega)$ is the macroscopic dielectric constant of the cell along $\qhat$.
It must be positive, $a$ being a quadratic form of the positive head $Y$ minus the body
correction; the code aborts if it is not.

\subsection{Cell integrals}
\label{sec:cell3d}

The Brillouin-zone average now carries $\Omega/(2\pi)^3$ and the volume element
$q^2\dd q\,\dd\Omega_{\qhat}$, which cancels the $1/q^2$ of the head, so that the radial
integral is elementary:
\begin{equation}
\begin{aligned}
  K_\head(\bR,\ii\omega) &= -\frac{\Omega}{(2\pi)^3}\int\dd\Omega_{\qhat}\,
  \frac{16\pi^2 a}{1+4\pi a}\int_0^{L(\qhat)}\dd q\,\cos(q\,\qhat\cdot\bR)\\
  &= -\frac{\Omega}{(2\pi)^3}\int\dd\Omega_{\qhat}\,\frac{16\pi^2 a}{1+4\pi a}\,
  \frac{\sin(Lt)}{t},
\end{aligned}
  \label{eq:Kh3d}
\end{equation}
with $t = \qhat\cdot\bR$ and $L(\qhat)$ the distance from $\Gamma$ to the boundary of the cell
$\{\sum_i s_i\mathbf b_i/2,\ |s_i|\le1\}$ along $\qhat$. In the wing the volume element combines
with the $1/q$ of Eq.~\eqref{eq:wing3d},
\begin{equation}
\begin{aligned}
  K_{\wing,\alpha}(\bR,\ii\omega) &= \frac{\Omega}{(2\pi)^3}\int\dd\Omega_{\qhat}\,
  \frac{4\pi\,\hat q_\alpha}{1+4\pi a}\int_0^{L(\qhat)}\dd q\,q\,\sin(q\,\qhat\cdot\bR)\\
  &= \frac{\Omega}{(2\pi)^3}\int\dd\Omega_{\qhat}\,\frac{4\pi\,\hat q_\alpha}{1+4\pi a}\,
  \frac{\sin(Lt) - Lt\cos(Lt)}{t^2}.
\end{aligned}
  \label{eq:Kw3d}
\end{equation}
Both radial results are Taylor expanded below $|Lt| = 0.02$, where the closed forms cancel,
\begin{equation}
\begin{aligned}
  \frac{\sin(Lt)}{t} &= L\Big(1 - \frac{(Lt)^2}{6} + \frac{(Lt)^4}{120} - \cdots\Big),\\
  \frac{\sin(Lt)-Lt\cos(Lt)}{t^2} &= L^2\Big(\frac{Lt}{3} - \frac{(Lt)^3}{30} + \frac{(Lt)^5}{840} - \cdots\Big),
\end{aligned}
  \label{eq:taylor}
\end{equation}
with a next omitted term of relative size $10^{-13}$.

The solid-angle integral is a cubature over the six faces of that cell. Face pair $i$ sits at
$s_i=\pm1$ and is parametrized by $(s_j,s_k)$, a point on it is
$\bq = \pm\mathbf b_i/2 + s_j\mathbf b_j/2 + s_k\mathbf b_k/2$, and its solid-angle element
seen from $\Gamma$ is
\begin{equation}
  \dd\Omega = \frac{|\bq\cdot\dd\mathbf S|}{q^3}
  = \frac{|\bq\cdot(\mathbf b_j\times\mathbf b_k)|}{4\,q^3}\,\dd s_j\,\dd s_k ,
  \label{eq:dOmega}
\end{equation}
integrated with an $8\times8$ Gauss--Legendre rule per face ($6\times64$ directions $\qhat_d$
with weights $\varpi_d$), so that
$\int_\mathrm{cell}\dd^3q\,f(\qhat)g(q) = \sum_d\varpi_d\int_0^{L(\qhat_d)}\dd q\,q^2f(\qhat_d)g(q)$.
Both signs of every face are included, which makes the direction set symmetric under
$\qhat\to-\qhat$ and the $\cos$/$\sin$ split of the kernels exact. The rule is checked by
$\sum_d \varpi_d L_d^3/3 = (2\pi)^3/\Omega$.

\subsection{Where the bulk correction is applied: in the orbital basis, not through the RI fit}
\label{sec:where3d}

In two dimensions the head is a $\sim30$~meV effect and Eq.~\eqref{eq:DeltaW} can be applied
in the auxiliary basis. In three dimensions the head potential $K_\head(0,\ii\omega)/\Omega$ is
of order eV ($-3.5$~eV for a $3\times3\times3$ diamond cell with $\varepsilon_\infty = 3$) and
scales as $1/L$, and the same application fails. $c_P c_Q K_\head(\bR_{PQ})$ is the projection
of a smooth, long-wavelength kernel onto atom-centred Gaussians; its fit through
$M^{-1}\,\cdot\,M^{-1}$ in Eq.~\eqref{eq:Wmic} demands large cancelling coefficients, and the
neglected form factors break the cancellation. On the 54-atom diamond cell the head alone,
applied in the RI basis, gives a gap of $-96$~eV and splits a degenerate triplet by $0.4$~eV
--- which a uniform potential cannot do. The bulk correction is therefore added where it is an
exact functional of the orbital-pair densities: to $\Sigma^c$ in the orbital basis, after
Eq.~\eqref{eq:Sc_ao}, with $\Delta W_{\ell\ell'}=0$ in Eqs.~\eqref{eq:Sc_neg}--\eqref{eq:Sc_pos}.
No RI fit of a kernel is involved anywhere.

\subsection{The orbital-pair kernel}
\label{sec:What}

Between orbital-pair densities $\rho_{nm} = \psi_n\psi_m$ the cell average of the correction
needs only the coupling of the pair to the charge mode,
\begin{equation}
  \langle mn|\phi_0\rangle = \frac{1}{\sqrt\Omega}\int\dd\br\,\rho_{mn}(\br)\,\ee^{\ii\bq\cdot\br}
  = \frac{\delta_{mn} + \ii\,\bq\cdot\mathbf d_{mn}}{\sqrt\Omega} + O(q^2),
  \label{eq:pair_phi0}
\end{equation}
with the velocity-gauge dipole between \emph{any} two orbitals,
\begin{equation}
  \mathbf d_{nm} = \langle n|\br|m\rangle
  = \frac{\langle n|[H,\br]|m\rangle}{\varepsilon_n - \varepsilon_m},
  \label{eq:dnm}
\end{equation}
which vanishes between folded states of different crystal momentum, where the commutator
vanishes, and is set to zero inside degenerate manifolds ($|\varepsilon_n-\varepsilon_m|<10^{-5}$~Ha),
where it is not defined. Collecting $\langle nm|\Delta W(\bq)|mn\rangle$ from the head
\eqref{eq:headwing3d} and the wings \eqref{eq:wing3d} and averaging over the cell gives the
real kernel
\begin{equation}
\begin{aligned}
  \hat W_{nm}(\ii\omega) ={}& w(\ii\omega)\,\delta_{nm}
  + \sum_{\alpha\beta} B^\hd_{\alpha\beta}(\ii\omega)\, d_{nm,\alpha} d_{nm,\beta}\\
  &+ 2\sqrt\Omega\sum_{\alpha\beta} A^\wing_{\alpha\beta}(\ii\omega)\,
  V_{nm,\alpha}(\ii\omega)\, d_{nm,\beta},
\end{aligned}
  \label{eq:What}
\end{equation}
with the three ingredients
\begin{align}
  w(\ii\omega) &= \frac{K_\head(\mathbf 0,\ii\omega)}{\Omega},
  \label{eq:headpot}\\
  A^\wing_{\alpha\beta}(\ii\omega) &= \frac{1}{(2\pi)^3}\int_\mathrm{cell}\dd^3q\;
  \frac{4\pi}{\tilde\varepsilon(\qhat,\ii\omega)}\,\hat q_\alpha\hat q_\beta,
  \label{eq:Awing}\\
  B^\hd_{\alpha\beta}(\ii\omega) &= \frac{1}{(2\pi)^3}\int_\mathrm{cell}\dd^3q\;
  4\pi\Big(\frac{1}{\tilde\varepsilon(\qhat,\ii\omega)} - 1\Big)\hat q_\alpha\hat q_\beta,
  \label{eq:Bhd}
\end{align}
evaluated with the face cubature as
$\int_\mathrm{cell}\dd^3q\,f(\qhat) = \sum_d\varpi_d\,L_d^3 f(\qhat_d)/3$, and the wing potential
\begin{equation}
\begin{aligned}
  V_{nm,\alpha}(\ii\omega) &= \int\dd\br\,\rho_{nm}(\br)\,w_\alpha(\br;\ii\omega)\\
  &= \sum_\ell \rho_{nm}(\br_\ell)\,\big[Z M^{-1} T(\ii\omega)\big]_{\ell\alpha},\\
  w_\alpha &= \sum_P \hat\phi_P\,T_{P\alpha},
\end{aligned}
  \label{eq:Vnm}
\end{equation}
where $\hat\phi_P$ denotes the RI-RS dual of $\phi_P$, i.e.\ the functional
$\int\rho\,\hat\phi_P = \sum_\ell\rho(\br_\ell)(ZM^{-1})_{\ell P}$ that maps a density to its
coefficient in the representation of $W^\MIC$. With $\tilde w = ZM^{-1}T$ this is evaluated as
\begin{equation}
  V_{nm,\alpha} = \big[C^\Tr\big(\Phi^\Tr\,\mathrm{diag}(\tilde w_{\cdot\alpha})\,\Phi\big)C\big]_{nm}.
  \label{eq:Vnm_grid}
\end{equation}
The head is a uniform potential (order $1/L$), the head-dipole term is the $O(q^2)$
remainder of the head, and the wing couples the dipole of the pair density to the potential
$w_\alpha$ that the uniform field induces in the body ($A^\wing, B^\hd\propto1/\Omega$ and
$T\propto1/\sqrt\Omega$, so both are of order $1/\Omega$). The wing carries a single
$1/\sqrt\Omega$ from Eq.~\eqref{eq:pair_phi0} against the $\Omega/(2\pi)^3$ of the average,
hence $2\sqrt\Omega$; the head terms carry $|\langle mn|\phi_0\rangle|^2\propto1/\Omega$, which
the average cancels. $V_{nm}$ is a linear functional of $\rho_{nm}$ and is exact on the RI-RS
grid, unlike the pair kernel $c_Pc_QK_\head(\bR_{PQ})$ of the RI-space application,
Eq.~\eqref{eq:DeltaW}.

\subsection{Correction of the correlation self-energy}
\label{sec:dSigma3d}

$\hat W_{nm}$ is transformed to imaginary time with the cosine weights of the minimax grid,
$\hat W_{nm}(\ii\tau_j) = \sum_k\tilde\gamma^{\cos}_{jk}\hat W_{nm}(\ii\omega_k)$. With the
$\Sigma^c$ construction of Eqs.~\eqref{eq:Sc_neg}--\eqref{eq:Sc_pos}
($\Sigma^c(-\ii|\tau|) = -G^\occ\circ W$, $\Sigma^c(+\ii|\tau|) = +G^\vir\circ W$) the
orbital-diagonal correction is
\begin{align}
  \Delta\Sigma^c_n(-\ii|\tau|) &= -\sum_{m}^{\occ} \ee^{-|\tau|(\varepsilon_F-\varepsilon_m)}\,\hat W_{nm}(\ii\tau),
  \label{eq:dSigma_neg}\\
  \Delta\Sigma^c_n(+\ii|\tau|) &= +\sum_{m}^{\vir} \ee^{-|\tau|(\varepsilon_m-\varepsilon_F)}\,\hat W_{nm}(\ii\tau),
  \label{eq:dSigma_pos}
\end{align}
added to the AO matrices of Eq.~\eqref{eq:Sc_ao} as
\begin{equation}
  \Delta\Sigma^c_{\mu\nu}(\pm\ii|\tau|) = \Big[(\Sov C)\,\mathrm{diag}\big(\Delta\Sigma^c_n(\pm\ii|\tau|)\big)\,(\Sov C)^\Tr\Big]_{\mu\nu},
  \label{eq:dSigma_ao}
\end{equation}
with $\Sov$ the AO overlap matrix. Because $C^\Tr\Sov C = 1$ this folds onto the band and
density-of-states $k$-points exactly like the regular $\Sigma^c$. For the head alone,
$\hat W_{nm} = w\,\delta_{nm}$, Eqs.~\eqref{eq:dSigma_neg}--\eqref{eq:dSigma_ao} reduce to
\begin{equation}
\begin{aligned}
  \Delta\Sigma^c_{\mu\nu}(-\ii|\tau|) &= -w(\ii\tau)\,\big[\Sov\,C_\occ\,\ee^{-|\tau|(\varepsilon_F-\varepsilon)}C_\occ^\Tr \Sov\big]_{\mu\nu},\\
  \Delta\Sigma^c_{\mu\nu}(+\ii|\tau|) &= +w(\ii\tau)\,\big[\Sov\,C_\vir\,\ee^{-|\tau|(\varepsilon-\varepsilon_F)}C_\vir^\Tr \Sov\big]_{\mu\nu},
\end{aligned}
  \label{eq:dSigma3d}
\end{equation}
whose orbital diagonal is $\mp w(\ii\tau)\,\ee^{-|\tau(\varepsilon_n-\varepsilon_F)|}$ for
occupied and empty $n$: the $\pm w/2$ shift of the COHSEX head in the static limit. A uniform
potential couples a state only to itself, $\langle m|\ee^{\ii\bq\cdot\br}|n\rangle =
\delta_{mn} + O(q)$, and therefore cannot split a degenerate manifold; the head-dipole and wing
terms carry the $O(\bq\cdot\mathbf d_{mn})$ corrections of the same expansion.

%======================================================================================
\section{Implementation in CP2K}
\label{sec:impl}
%======================================================================================

\subsection{Activation}
\label{sec:activation}

The correction lives in the module \texttt{gw\_gamma\_correction} and is switched on
automatically for the large-cell $\Gamma$ drivers, the RI-RS one and the exact-RI tensor one
(\texttt{gw\_tensor\_large\_cell\_gamma}), when (i) the system is a slab
(\texttt{PERIODIC XY}, head $2\pi/q$) or a bulk cell (\texttt{PERIODIC XYZ}, head $4\pi/q^2$),
and (ii) the mesh for $\chi$, $\varepsilon$, $W$ is \texttt{KPOINTS\_W 1 1 1}. Only the
$1\times1\times1$ mesh samples $\bq=0$: the setup keeps $n_1n_2n_3/2$ points of an
$n_1\times n_2\times n_3$ mesh, one per $\pm\bk$ pair, and for odd $n>1$ the integer division
drops the self-paired $\Gamma$ point. For this case the extrapolation mesh is dropped (otherwise
$\Gamma$ would enter the $s$--$s$ elements of $W$ with the extrapolation weight $-2.414$ and
every other element with weight 0). Wires and slabs whose vacuum direction is not $z$ are not
implemented. There is no input keyword. In the tensor driver only the optical response
$S_{\alpha\beta}$, $T_{P\alpha}$ of Sec.~\ref{sec:charge} is evaluated on the RI-RS grid, which
is built for that purpose; $\chi$, the body of $W$ and $\Sigma^c$ remain exact-RI tensor
contractions. This gives a reference in which the RI-RS fit enters only through the
correction, and it is used in Sec.~\ref{sec:results} to separate fit errors from the
$\Gamma$-point model.

\subsection{Call sequence}
\label{sec:calls}

Table~\ref{tab:calls} lists the routines in the order in which the $GW$ driver calls them.
\texttt{gamma\_corr\_init} runs before the $Z_{\ell P}$ fit and sets only geometric quantities
(layer area or cell volume, reciprocal vectors, metric head). The Coulomb side, the response and
the cell integrals are built in step 4b of the $\Gamma$ driver
(\texttt{gw\_ri\_rs\_large\_cell\_gamma} or \texttt{gw\_tensor\_large\_cell\_gamma}), after
$\chi_{PQ}(\Gamma,\ii\tau)$ exists and before \texttt{get\_W\_MIC} consumes it. The $V(0)$ of
\texttt{gamma\_corr\_coulomb\_at\_zero} is kept distributed and reused by the regular $W$ path
instead of repeating the lattice sum. In 2D, \texttt{gamma\_corr\_add\_to\_W} adds
$\Delta W_{PQ}(\ii\tau)$ inside \texttt{compute\_W\_MIC}, before the contraction
$M^{-1}W^\MIC M^{-1}$; in 3D it returns immediately and
\texttt{gamma\_corr\_add\_to\_sigma} adds Eqs.~\eqref{eq:dSigma_neg}--\eqref{eq:dSigma_ao}
after \texttt{compute\_Sigma\_c} and before the quasiparticle solve.

\begin{table*}
\caption{Call sequence of the $\Gamma$-point correction in \cpk{}. ``2D''/``3D'' mark
dimension-specific branches.}
\label{tab:calls}
\footnotesize
\begin{ruledtabular}
\begin{tabular}{>{\raggedright}p{0.25\textwidth}>{\raggedright}p{0.30\textwidth}>{\raggedright\arraybackslash}p{0.38\textwidth}}
caller & routine & computes \\
\colrule
\texttt{ri\_rs\_setup} & \texttt{gamma\_corr\_init} & activation; $A$ or $\Omega$, $\mathbf b_i$, metric head; dense-algebra mode\\
$\Gamma$ driver, step 4b & \texttt{gamma\_corr\_build\_coulomb\_maps} & $c$, $\Pi$, $V_b$, $F^\dens$, $G^\ext$; $V(0)$ via \texttt{gamma\_corr\_coulomb\_at\_zero}\\
 & \texttt{gamma\_corr\_response} & dipoles, $Y$, $X$, $B$, $E$; $S_{\alpha\beta}(\ii\omega)$, $T_{P\alpha}(\ii\omega)$\\
 & \texttt{gamma\_corr\_tabulate} & 2D: $K_\head(\bR)$, $K_\wing(\bR)$ tables; 3D: \texttt{gamma\_corr\_direction\_coefficients}\\
\texttt{gw\_compute\_V\_and\_W} & \texttt{gamma\_corr\_subtract\_divergent\_head} & 2D only: $V(0)-D_\mathrm{box}cc^\Tr$, Eq.~\eqref{eq:Vreg}\\
\texttt{compute\_W\_MIC} & \texttt{gamma\_corr\_add\_to\_W} & 2D: $W^\MIC_{PQ}(\ii\tau) \mathrel{+}= \Delta W_{PQ}(\ii\tau)$, Eq.~\eqref{eq:DeltaW}\\
$\Gamma$ driver, after $\Sigma^c$ & \texttt{gamma\_corr\_add\_to\_sigma} & 3D: $\Sigma^c_{\mu\nu}(\pm\ii|\tau|) \mathrel{+}= \Delta\Sigma^c_{\mu\nu}$, Eq.~\eqref{eq:dSigma_ao}\\
\end{tabular}
\end{ruledtabular}
\end{table*}

The wing quantity $X_{\ell\alpha}$ of Eq.~\eqref{eq:X} is evaluated as the diagonal of
$\Phi P^\alpha(\ii\tau)\Phi^\Tr$ with the dense AO matrix
$P^\alpha = C_\occ\,\ee\,\mu^\alpha\,\ee\,C_\vir^\Tr$, using that the exponential factorizes,
$\ee^{-(\varepsilon_a-\varepsilon_i)\tau} = \ee^{-|\varepsilon_a-\varepsilon_F|\tau}\ee^{-|\varepsilon_i-\varepsilon_F|\tau}$;
this is formed by two ScaLAPACK multiplications and one sparse product on the block pattern
of $\Phi$, and is correct under any distribution of the MO coefficients.

\subsection{Replicated or distributed dense algebra}
\label{sec:dense}

The Coulomb maps and the response involve dense $N_\mathrm{RI}\times N_\mathrm{RI}$ algebra
($V$, $\Pi$, $V_b$, $V_b^{\pm1/2}$, $M^{-1}$, $F^\dens$, $G^\ext$, $B$ and its inverse). The
replicated path holds up to ten such arrays per rank. \texttt{gamma\_corr\_init} estimates that
peak as $10\times8N_\mathrm{RI}^2$~bytes, measures the memory available per MPI rank, and takes
the distributed (ScaLAPACK) path when the estimate exceeds a quarter of it; the remaining
three quarters are left for the distributed matrices of the regular $W$ path. With 12 auxiliary
functions per atom the replicated peak is 4~GB per rank at about 600 atoms and 46~GB at 2000
atoms. In the distributed path $\Pi$ is filled locally, $V_b = \Pi V\Pi$, $F^\dens$ and $G^\ext$
are parallel matrix products, $V_b^{\pm1/2}$ come from one distributed eigendecomposition with
the same relative eigenvalue cut, and per frequency $B$ is inverted by a Cholesky factorization,
$x = B^{-1}E$, $S \mathrel{-}= E^\Tr x$, $T = G^\ext x$. Only $N_\mathrm{RI}\times3$ arrays are
replicated. The choice is automatic; the two paths agree to all printed digits on one rank.
The remaining replicated arrays, the dipoles $\mu$ ($3\,n_\occ n_\vir$) and $X(\ii\tau)$ on the
grid ($3N_\mathrm{grid}N_\tau$), grow linearly with system size.

\subsection{Cost}
\label{sec:cost}

In 2D the correction costs one table lookup per auxiliary pair and time point in
\texttt{gamma\_corr\_add\_to\_W}; the tables are built once from the $3\times3$ tensors $S$ on a
$128\times128$ mesh. In 3D the cell integrals are per-direction coefficients
($6\times8\times8$ directions) and the radial integrals are analytic. The orbital-pair kernel
costs, per frequency, three products $\Phi^\Tr\mathrm{diag}(\tilde w_{\cdot\alpha})\Phi$ on the
grid, six $n_\mathrm{AO}^2n_\mathrm{MO}$ transforms to the MO basis, and an
$n_\mathrm{MO}^2/n_\mathrm{ranks}$ element loop; the dipoles $\mathbf d_{nm}$ need three
$n_\mathrm{AO}^2n_\mathrm{MO}$ transforms once per spin.
% TIMINGS TO BE FILLED BY RITAJ (2D and 3D wall times of the correction)

\subsection{Order of evaluation}
\label{sec:order}

\begin{enumerate}
\item $\Phi_\mu(\br_\ell)$, Eq.~\eqref{eq:phi_grid}; $Z_{\ell P}$ sphere by sphere, Eqs.~\eqref{eq:D}--\eqref{eq:fit_reg}.
\item $G^\occ,G^\vir$, $\chi_{\ell\ell'}$, $\chi_{PQ}(\ii\tau)\to\chi_{PQ}(\ii\omega)$, Eqs.~\eqref{eq:G_ao}--\eqref{eq:chi}.
\item $c_P$, $\Pi$, $V_b$, $F^\dens$, $F^\ext$, Eqs.~\eqref{eq:charge}--\eqref{eq:F}; dipoles
  and sum rule, Eqs.~\eqref{eq:dipole}--\eqref{eq:TRK}; $Y$, $X$, $B$, $E$, $S$, $T$,
  Eqs.~\eqref{eq:Y}--\eqref{eq:ST}; 2D: tables $K_\head$, $K_\wing$,
  Eqs.~\eqref{eq:Kh}--\eqref{eq:Kw}; 3D: direction coefficients, $w$, $A^\wing$, $B^\hd$,
  Eqs.~\eqref{eq:Kh3d}, \eqref{eq:headpot}--\eqref{eq:Bhd}.
\item 2D: $V^\lat(0)\to V(0)$, Eqs.~\eqref{eq:Vlat}--\eqref{eq:Vreg}; 3D: the Ewald $V(0)$ of
  Eq.~\eqref{eq:Vewald}, nothing subtracted. $\varepsilon$, $W^c$, $W^c(\ii\omega)\to W^c(\ii\tau)$;
  2D: $W^c_{PQ}(\ii\tau) \mathrel{+}= \Delta W_{PQ}(\ii\tau)$, Eq.~\eqref{eq:DeltaW}; then
  $W^\MIC$, Eqs.~\eqref{eq:eps}--\eqref{eq:Wmic}.
\item $\Sigma^x$, Eq.~\eqref{eq:Sx}.
\item Per time point and panel: $W^\MIC_{\ell\ell'}$, Eq.~\eqref{eq:W_grid}; $\Sigma^c$ in both
  branches, Eqs.~\eqref{eq:Sc_neg}--\eqref{eq:Sc_ao}.
\item 3D: $\hat W_{nm}$, Eqs.~\eqref{eq:dnm}--\eqref{eq:Vnm_grid}; $\Delta\Sigma^c$,
  Eqs.~\eqref{eq:dSigma_neg}--\eqref{eq:dSigma_ao}.
\item Eqs.~\eqref{eq:Sc_freq}--\eqref{eq:QP}.
\end{enumerate}

%======================================================================================
\section{Results}
\label{sec:results}
%======================================================================================

All reference calculations use the full-$k$-point RI-RS $GW$ of J.~Wilhelm, cp2k
PR~\#6447, in the form in which it was merged into the \cpk{} master branch (commit
\texttt{7b88fc63f1}). The $\Gamma$-only code was rebased onto the same master, so that every
number of this section, ours and the reference, comes from one binary. That master also
carries a change of the analytic continuation (PR~\#6474: the quasiparticle energy is the
average over several central Pad\'e orders instead of a single 16-parameter fit), which
shifts every $G_0W_0$ gap of this section by $+0.2$ to $+0.4$~eV relative to our earlier
standard-Pad\'e numbers, identically in both codes, and leaves Hartree--Fock energies
untouched; the differences between the two codes are unchanged within 6~meV. Both codes use
the same orbital and auxiliary bases, pseudopotentials, cutoffs, minimax grid and RI
settings. The PR requires its own optimised RI-RS grid
(\texttt{RS\_AO\_RATIO}~7, about 90 points per atom with a DZVP basis); a three times denser
grid changes the h-BN gap by 2~meV. Our runs use the tabulated atom-centred grids of CP2K
(\texttt{data/ri\_rs\_grid}, about 340 points per atom). Odd meshes of the PR abort in its own
$\Gamma$ treatment, so its series are $12$ to $36$ in steps of 6 (2D) and $4$ to $12$ in
steps of 2 (3D).

\subsection{Validation of the slab correction: h-BN with the aug-SZV basis}
\label{sec:resval}

Table~\ref{tab:resval} compares the $\Gamma$-only $8\times8$ h-BN supercell (128 atoms,
aug-SZV-MOLOPT-GTH tier-1 orbital and RI bases, 15~\AA\ vacuum, cutoff 600~Ry, 10 minimax
points) with the PR at increasing meshes, at the primitive point $P=(11/32,11/32)$ nearest to
$K$ that the supercell samples, and state by state at the point $Q$ that the supercell folds
onto its $\Gamma$. The full-$k$ gap converges from below towards the $\Gamma$-only value; at
$24\times24$ the absolute quasiparticle energies at $Q$ agree to 16~meV (valence) and 24~meV
(conduction), with identical $\Sigma^x$ and $v^{xc}$, and the gap at $P$ to 56~meV. The larger
$\Gamma$-only cells sample a different point nearest to $K$ ($13/40$ and $19/56$ instead of
$11/32$) and move the gap by 11 and 34~meV; the full-$k$ gap itself differs by 37~meV between
$K$ and $P$, so these are sampling shifts rather than cell convergence. With the standard
Pad\'e continuation of our earlier builds the same comparison gave 5 and 14~meV at $Q$ and
20~meV at $P$ (Table~\ref{tab:oldscheme}); the averaged continuation amplifies the small
$\Sigma^c$ differences between the folded supercell spectrum and the primitive-cell spectrum
(2.043 vs 2.082~eV at the valence state of $Q$ at the same 8.7~\AA\ exchange cutoff).

\begin{table}[h]
\caption{h-BN monolayer, aug-SZV tier-1 basis: direct $G_0W_0$ gap at $P$ and quasiparticle
energies at $Q$ (eV); one master build, averaged Pad\'e continuation in both codes.
$\Gamma$-only: $8\times8$ supercell with the head-and-wing correction (the $10\times10$ and
$14\times14$ entries are the gap at the sampled point nearest to $K$ of those cells, $13/40$
and $19/56$). PR: primitive cell, $N\times N$ mesh.}
\label{tab:resval}
\begin{ruledtabular}
\begin{tabular}{lcccc}
 & gap($P$) & VB($Q$) & CB($Q$) & $\Sigma^c$ VB / CB\\
\colrule
PR $6\times6$   & 6.714 & $-5.013$ & 2.169 & 2.050 / $-3.088$\\
PR $8\times8$   & 6.780 & $-5.051$ & 2.161 & 2.082 / $-3.113$\\
PR $12\times12$ & 6.801 & $-5.082$ & 2.143 & 2.118 / $-3.146$\\
PR $16\times16$ & 6.804 & $-5.096$ & 2.129 & 2.134 / $-3.161$\\
PR $20\times20$ & 6.807 & $-5.105$ & 2.122 & 2.142 / $-3.168$\\
PR $24\times24$ & 6.810 & $-5.112$ & 2.119 & 2.146 / $-3.172$\\
$\Gamma$-only $8\times8$ & 6.866 & $-5.128$ & 2.143 & 2.043 / $-3.148$\\
$\Gamma$-only $10\times10$ & 6.855 & & & \\
$\Gamma$-only $14\times14$ & 6.832 & & & \\
\end{tabular}
\end{ruledtabular}
\end{table}

\subsection{2D benchmark with a DZVP basis}
\label{sec:res2d}

Four planar hexagonal monolayers (h-BN, SiC, h-GaN, h-AlN) with DZVP-MOLOPT-SR-GTH orbital
bases and MOLOPT RI sets at the $2\cdot10^{-3}$ fit-error level, 15~\AA\ vacuum, cutoff
500~Ry, 10 minimax points, \texttt{CUTOFF\_RADIUS\_RI} 5~\AA, \texttt{SIZE\_LATTICE\_SUM}~4.
The $\Gamma$-only cells are $9\times9$ (162 atoms; the folded manifold contains $K$) and
$12\times12$ (288 atoms); the PR meshes contain $K$ as well. Table~\ref{tab:res2d} lists the
indirect gap over the sampled points. The PR series is converged to 5~meV between
$30\times30$ and $36\times36$, the $\Gamma$-only results to 3--15~meV between $9\times9$ and
$12\times12$. Three of the four materials agree with the converged reference to
$22$--$74$~meV, and their Hartree--Fock gaps to $0.01$--$0.05$~eV (h-BN 13.267 vs 13.278, SiC
8.399 vs 8.373, h-GaN 9.733 vs 9.779~eV). h-AlN deviates by $0.28$~eV in the $G_0W_0$ gap and
by $0.13$~eV in the Hartree--Fock gap. State by state ($9\times9$ against PR $36\times36$)
the difference is spread over four contributions of the same sign: $\Sigma^x$ at the valence
maximum at $K$ ($-19.081$ vs $-19.149$~eV) and at the nearly-free-electron conduction minimum
at $\Gamma$ ($-4.894$ vs $-4.834$~eV), together $0.13$~eV, and $\Sigma^c$ at the valence
maximum ($2.463$ vs $2.414$~eV) and at the conduction minimum ($-2.448$ vs $-2.347$~eV),
together $0.15$~eV; $\varepsilon^\mathrm{DFT}$ and $v^{xc}$ of the four states are identical
to the meV. Neither grid is responsible. The PR grid is converged for AlN
(\texttt{RS\_AO\_RATIO} 20 gives 6.318~eV against 6.319~eV at the default ratio 7), and our
$\Sigma^x$ is converged with the truncated-Coulomb radius (12.1~\AA\ at $9\times9$, 16.2~\AA\
at $12\times12$: 4~meV). On our side the exact-RI tensor driver with the same $\Gamma$-only $W$
and the same correction (Sec.~\ref{sec:activation}) reproduces the RI-RS result of the
$6\times6$ AlN cell to 1~meV in every occupied and low-lying virtual state: gap 6.241~eV in
both, and $\varepsilon^\mathrm{DFT}$, $\Sigma^c$, $\Sigma^x$, $v^{xc}$ of the two band edges identical
to the printed meV ($-19.063$ and $-15.173$~eV at the valence maximum, $-4.901$ and
$-8.832$~eV at the conduction minimum for $\Sigma^x$ and $v^{xc}$), at ten times the cost per
$\chi$ time point (33 vs 3~s) and six times the memory. The remaining AlN difference therefore lies in the
$\Gamma$-point treatment itself, or in the supercell-$\Gamma$ versus full-$k$ formulation
(32--56~\AA\ Coulomb truncation in the PR against 12--16~\AA\ here, folded spectrum on the
minimax grid here), and not in
the numerics of either code; see Sec.~\ref{sec:disc_approx}.

\begin{table}[h]
\caption{2D benchmark, DZVP basis ($a$ = 2.502, 3.090, 3.210, 3.110~\AA): indirect $G_0W_0$ gap
(eV) of the $\Gamma$-only supercell with the head-and-wing correction against the full-$k$ PR
at $N\times N$ meshes; one master build, averaged Pad\'e continuation in both codes. Last
column: $\Gamma$-only ($9\times9$) minus PR $36\times36$.}
\label{tab:res2d}
\begin{ruledtabular}
\footnotesize
\begin{tabular}{lcccccccc}
system & $9\times9$ & $12\times12$ & PR 12 & PR 18 & PR 24 & PR 30 & PR 36 & $\Delta$\\
\colrule
h-BN  & 7.837 & 7.834 & 7.827 & 7.839 & 7.847 & 7.854 & 7.859 & $-0.022$\\
SiC   & 4.224 & 4.221 & 4.261 & 4.270 & 4.273 & 4.275 & 4.278 & $-0.054$\\
h-GaN & 4.464 & 4.479 & 4.514 & 4.527 & 4.532 & 4.535 & 4.538 & $-0.074$\\
h-AlN & 6.299 & 6.303 & 6.541 & 6.555 & 6.565 & 6.571 & 6.576 & $-0.277$\\
\end{tabular}
\end{ruledtabular}
\end{table}

\subsection{3D: diamond}
\label{sec:res3d}

Diamond ($a=3.567$~\AA, primitive two-atom cell, SZV-MOLOPT-SR-GTH, RI set
\texttt{RI\_20\_2\_2\_1\_1\_0\_0\_0}, cutoff 400~Ry, 10 minimax points,
\texttt{SIZE\_LATTICE\_SUM}~2, RI-metric cutoff 3.0~\AA, $\Sigma^x$ truncated at half the
cell height) in $n\times n\times n$ supercells of 54, 128 and 250 atoms (the 432-atom cell was
abandoned after seven hours on four nodes, still in the fit stage). Table~\ref{tab:res3d} shows the
indirect $G_0W_0$ gap over the folded $\Gamma$ manifold for the successive stages of the
bulk correction, together with the PR at $N\times N\times N$ meshes and the Hartree--Fock
gaps. Two things are established independently of the PR. First, the exchange: the exact-RI
tensor implementation of the same 128-atom cell gives $\Sigma^x=-18.957$~eV (valence triplet)
and $-9.767$~eV (conduction) and a Hartree--Fock gap of 15.125~eV, identical to the RI-RS
values ($-18.957$/$-9.767$, 15.123~eV), whereas the PR at $4\times4\times4$ is 0.4--0.6~eV off
in $\Sigma^x$ and splits the $\Gamma$ triplet by 0.2~eV, reaching the correct exchange only at
$8\times8\times8$. Second, the structure of the correction: the head alone shifts the valence
and conduction edges by $+1.58$ and $-1.80$~eV at $3\times3\times3$ (the $\pm w/2$ of a uniform
potential with $w=-3.5$~eV, $\varepsilon_\infty=3.05$); the head-dipole term adds $+0.20$~eV
to the gap at 54 atoms and $+0.06$~eV at 128 atoms, the wing $-0.03$ and $-0.00$~eV, and the
wing-wing term $-0.002$~eV, i.e.\ the dipole terms scale as $1/\Omega$ as they should. The
RI-space application of the same head, by contrast, gave $-96$~eV. Degeneracies of the folded
manifold are preserved to 3~meV. With the full kernel and the one master build the
$\Gamma$-only gap runs $9.210\to9.970\to10.100$~eV (54, 128, 250 atoms; the last step
$+0.13$~eV), the PR $9.239\to9.519\to9.593\to9.553\to9.565$~eV ($4^3$ to $12^3$ meshes),
converged to 0.01~eV from $8^3$ on. The Hartree--Fock gaps rise on both sides towards the same
value, $15.216$~eV at 250 atoms and $15.120$~eV at $12^3$ with $+0.04$~eV per step left in the
PR, so the exchange is consistent to 0.1~eV while the $G_0W_0$ gap of the $\Gamma$-only cell
lies $0.5$~eV above the converged full-$k$ value at 250 atoms. The whole of that difference is
in $\Sigma^c$.

Two further checks bound the numerics on our side. First, the exact-RI tensor driver with the
same $\Gamma$-only $W$ and the same orbital-basis correction (Sec.~\ref{sec:activation}) was
run on the 54- and 128-atom cells. On the 128-atom cell it gives the same gap as the RI-RS
code, 9.970~eV, and $\varepsilon^\mathrm{DFT}$, $\Sigma^c$, $\Sigma^x$, $v^{xc}$ of every
state of the valence triplet and of the conduction minimum identical to the printed meV
($1.950/1.932/1.932$, $-18.957$, $-17.093$~eV and $-3.206$, $-9.767$, $-16.568$~eV). On the
54-atom cell (standard-Pad\'e builds) its gap is 9.386~eV against 9.039~eV from the RI-RS
grid, but the difference is almost entirely exchange: the Hartree--Fock gaps are 15.074 and 14.838~eV, the
RI-RS $\Sigma^x$ is too negative by 0.06~eV on average over the occupied states and by 0.14~eV
over the lowest virtuals, whereas $\Sigma^c$ at the two band edges agrees to 0.004 and
0.036~eV. The $3\times3\times3$ cell truncates the Coulomb operator of $\Sigma^x$ at 3.1~\AA,
the same length as the 3.0~\AA\ RI-metric cutoff that defines the fit, and the grid
representation of such a short truncated operator is where the RI-RS fit is weakest; at
$4\times4\times4$ (4.1~\AA) the RI-RS and exact-RI $\Sigma^x$ agree to 1~meV, as stated
above. The 54-atom RI-RS entries of Table~\ref{tab:res3d} are therefore grid-limited by
$\sim0.2$~eV, the larger cells are not, and the $\Gamma$-only result of the 128-atom cell,
including the correction, is independent of the RI-RS fit. Second, the standard tensor $k$-point $GW$ of CP2K on the primitive
cell was attempted as an independent arbiter; with the 9~\AA\ orbital and RI radii its
real-space lattice sums (1607 cells of three-centre integrals, 12071 cells $\Delta\bR$)
cost 10.5 and 21.7~h per imaginary-time point on 8 ranks for the $6^3$ and $8^3$ meshes,
the code has no $\chi$ restart, and the runs were abandoned at the 24~h limit. The 0.5~eV
between the two $\Gamma$-point treatments at convergence is thus not settled by a third
code; what is settled is that it is not a fit error of either.

\begin{table}[h]
\caption{Diamond, SZV basis: indirect $G_0W_0$ gap over the folded $\Gamma$ manifold (eV) for
the stages of the bulk correction (columns none to full: development builds with the standard
Pad\'e continuation), the full kernel with the one master build and the averaged Pad\'e
continuation (full$^\ast$), and the full-$k$ PR from the same master build. HF: Hartree--Fock
gap with the truncated Coulomb operator of each run (master build).}
\label{tab:res3d}
\begin{ruledtabular}
\footnotesize
\begin{tabular}{lccccccc}
cell & atoms & none & head & +h.-dip. & full & full$^\ast$ & HF\\
\colrule
$3^3$ & 54  & 12.249 & 8.871 & 9.066 & 9.041 & 9.210 & 14.833\\
$3^3$, exact RI & 54 & & & & 9.386 & & 15.074\\
$4^3$ & 128 &        & 9.722 &       & 9.782 & 9.970 & 15.123\\
$4^3$, exact RI & 128 & & & & & 9.970 & 15.125\\
$5^3$ & 250 &        & 9.910 &       & 9.934 & 10.100 & 15.216\\
\colrule
PR $4^3$    & & & & & & 9.239 & 14.116\\
PR $6^3$    & & & & & & 9.519 & 14.784\\
PR $8^3$    & & & & & & 9.593 & 15.011\\
PR $10^3$   & & & & & & 9.553 & 15.077\\
PR $12^3$   & & & & & & 9.565 & 15.120\\
\end{tabular}
\end{ruledtabular}
\end{table}

%======================================================================================
\section{Discussion}
\label{sec:discussion}
%======================================================================================

\subsection{Why the old k-mesh large-cell scheme is not a reference}
\label{sec:oldscheme}

The large-cell $\Gamma$-point $GW$ implementations in \cpk{}
(\texttt{tensor\_large\_cell\_Gamma}, \texttt{ri\_rs\_large\_cell\_Gamma}) build the screened
interaction of a 2D supercell on a $k$-mesh \texttt{KPOINTS\_W}~$n\,n\,1$ of the supercell's
small Brillouin zone, with a second ``extrapolation'' mesh of $2n$ points and a Coulomb matrix
$V(\bq)$ from a minimum-image truncated real-space lattice sum. Band gaps from this scheme were
used as the reference against which the $\Gamma$-only scheme was judged, and the two differ by
$0.12$~eV on the largest cell we can afford. The old scheme carries the systematic error; the
error is the spurious net charge of RI-fitted pair densities identified by J.~Wilhelm in
\cpk{} PR~\#6447~\cite{pr6447}; and the $\Gamma$-only numbers are confirmed by his
full-$k$-point code to 15--60~meV (Sec.~\ref{sec:resval}). Numbers that came out of the old scheme for 2D systems
should therefore not be used as references.

\paragraph{The observation.} Monolayer h-BN, $a = 2.502$~\AA, 15~\AA\ vacuum, PBE,
aug-SZV-MOLOPT-GTH tier-1 orbital and RI bases, 10 minimax points,
\texttt{CUTOFF\_RADIUS\_RI}~5~\AA, identical in every code and run. ``Gap'' is the direct
$G_0W_0$ gap at the $k$-point $P = (11/32, 11/32)$ of the primitive zone, the point nearest to
$K$ that an $8\times8$ supercell with a $2\times2$ DOS mesh samples (Table~\ref{tab:oldscheme}).

\begin{table}[h]
\caption{Direct $G_0W_0$ gap of monolayer h-BN at $P=(11/32,11/32)$ from the old large-cell
scheme, the $\Gamma$-only scheme of this work, and the full-$k$ RI-RS $GW$ of
PR~\#6447 (primitive cell). All entries of this table are from the development builds with
the standard Pad\'e continuation, which is the only way to compare with the old scheme; the
one-build numbers are in Table~\ref{tab:resval}.}
\label{tab:oldscheme}
\begin{ruledtabular}
\begin{tabular}{llr}
scheme & cell / mesh & gap (eV)\\
\colrule
old, \texttt{KPOINTS\_W 4 4 1} + extrap. & $8\times8$ supercell & 8.152\\
old, \texttt{KPOINTS\_W 4 4 1} + extrap. & $14\times14$ supercell & 6.582\\
old, \texttt{KPOINTS\_W 2 2 1} + extrap. & $14\times14$ supercell & 6.595\\
\colrule
$\Gamma$-only + head/wing (this work) & $8\times8$ supercell & 6.710\\
$\Gamma$-only + head/wing & $10\times10$ supercell & 6.705\\
$\Gamma$-only + head/wing & $14\times14$ supercell & 6.699\\
\colrule
full-$k$ RI-RS, PR \#6447 & $6\times6$ mesh & 6.590\\
full-$k$ RI-RS, PR \#6447 & $8\times8$ mesh & 6.658\\
full-$k$ RI-RS, PR \#6447 & $12\times12$ mesh & 6.681\\
full-$k$ RI-RS, PR \#6447 & $16\times16$ mesh & 6.684\\
full-$k$ RI-RS, PR \#6447 & $20\times20$ mesh & 6.687\\
full-$k$ RI-RS, PR \#6447 & $24\times24$ mesh & 6.690\\
\end{tabular}
\end{ruledtabular}
\end{table}

Three things stand out. (i) The old scheme is not converged with the cell: $8.15$~eV at
$8\times8$ against $6.58$~eV at $14\times14$, and an earlier sweep with the tier-2 basis
(11.85, 9.74, 9.66, 9.43, 8.63, 7.54~eV for $n = 5, 6, 7, 8, 10, 12$) was still falling at
30~\AA. (ii) The $\Gamma$-only scheme is converged to 5~meV in the gap from $8\times8$ on.
(iii) The independent full-$k$-point code converges from below onto the $\Gamma$-only value;
state by state at the supercell $\Gamma$ point the absolute $G_0W_0$ energies agree to 5~meV
(valence) and 14~meV (conduction) with the $24\times24$ mesh, with identical $\Sigma^x$ and
$v^{xc}$ and $\Sigma^c$ within 10~meV. The old scheme sits $0.12$~eV below both at
$14\times14$ and is still drifting downwards, so the $0.12$~eV is its error, not the new
code's.

\paragraph{The mechanism: a spurious charge in the fitted pair densities.}
In RI-RS a product of two orbitals is interpolated from its values on the atom-centred grid
points $\br_\ell$ and fitted to the auxiliary basis through the fit coefficients $Z_{\ell P}$,
\begin{equation}
\begin{aligned}
  \psi_i(\br)\psi_a(\br) &\approx \sum_{P}\Big[\sum_\ell Z_{\ell P}\,\psi_i(\br_\ell)\psi_a(\br_\ell)\Big]\phi_P(\br)\\
  &\equiv \sum_P d_{ia,P}\,\phi_P(\br).
\end{aligned}
\end{equation}
The exact transition density is neutral, $\int\psi_i\psi_a = 0$ by orthogonality. The fitted one
is not: its charge $m_{ia} = \sum_P c_P\,d_{ia,P}$ is a small but nonzero number set by the grid
and the fit (the grid cannot represent orthogonality, and the fit is regularised). This was
found by J.~Wilhelm while converging his full-$k$-point RI-RS $GW$.

A spurious charge is harmless for a molecule and nearly harmless for a 3D crystal, but not in
two dimensions. On the $\bq$ mesh of a 2D calculation the Coulomb kernel in the charge channel
is $v(\bq) = 2\pi/(Aq)$, so the charge mode of $\chi$ enters the dielectric function as
\begin{equation}
  \varepsilon_{00}(\bq) - 1 \;\sim\; \frac{2\pi}{Aq}\sum_{ia}\frac{2\,|m_{ia}|^2\,\Delta_{ia}}{\Delta_{ia}^2+\omega^2}
  \;=\; \frac{C}{q},
\end{equation}
a $1/q$ divergence with a spurious prefactor. Summed over a 2D mesh the mesh points approach
$\bq = 0$ as $1/N$, and $\sum_\bq 1/q$ over an $N\times N$ mesh grows like $\ln N$: the
screening, and with it the gap, drifts logarithmically with the mesh instead of converging. In
Wilhelm's calculations the WSe$_2$ gap moved from 2.222 to 2.197~eV with the mesh; after
projecting the charge out of every fitted density, $d' = d - c\,(c^\dagger d)/(c^\dagger c)$, it
is 2.233~eV and stable; the 2D h-BN regression reference moved by $+0.92$~eV.

The old large-cell scheme has the same fitted densities with the same spurious charge, but its
Coulomb matrix is not $2\pi/(Aq)$. $V(\bq)$ is a real-space lattice sum truncated at the
minimum-image radius $R_c \approx L/2$ of the supercell, whose charge-channel value is finite,
$\int_0^{R_c}\!2\pi r\,\dd r/r = 2\pi R_c$ per cell area. The leak therefore does not
diverge; it produces a finite spurious screening proportional to $|m|^2 R_c/A \propto 1/L$,
i.e.\ a uniform downward shift of $\Sigma^c$ that decays only as the inverse cell size. That is
exactly the signature in Table~\ref{tab:oldscheme}: a gap that is too small, converges slowly
from below, and is still moving at 30~\AA. The sign of Wilhelm's fix ($+0.92$~eV for h-BN with
his settings) and the sign of our discrepancy (old scheme below) agree.

\paragraph{Why the $\Gamma$-only scheme is immune.}
The $\Gamma$-only implementation never lets a fitted density interact through the charge
channel of a sampled $V(\bq)$:
\begin{itemize}
\item the body of the dielectric matrix of the correction is built with the neutral projector,
      $V_b = \Pi V(0)\Pi$, Eq.~\eqref{eq:Pi}, so the charge direction of the auxiliary basis
      carries no Coulomb interaction at all in that body;
\item the head and wings of $\varepsilon$ in the charge channel are built from the exact
      velocity-gauge transition dipoles, Eq.~\eqref{eq:dipole}, checked by the sum rule of
      Eq.~\eqref{eq:TRK}, and integrated analytically over the whole small Brillouin zone with
      the $2\pi/q$ kernel;
\item the regular $W(\Gamma)$ uses $V(0)$ with only the divergent $\bG = 0$ constant removed,
      Eq.~\eqref{eq:Vreg}, and the charge direction of that matrix is quenched in
      $\varepsilon^{-1}$.
\end{itemize}
This is the same treatment of the charge channel as Wilhelm's $\Gamma$ model, which is why the
two codes agree once his leak is removed. The agreement (15~meV at $8\times8$ versus
$24\times24$, with the remaining difference of the size expected from the $8\times8$ cell) is
the actual validation of the $\Gamma$-only scheme; the old large-cell numbers are not.

\paragraph{Consequences.} Band gaps of 2D systems produced by
\texttt{tensor\_large\_cell\_Gamma} or \texttt{ri\_rs\_large\_cell\_Gamma} with
\texttt{KPOINTS\_W} $n\,n\,1$ ($n>1$) carry a cell-size-dependent error of order $0.1$~eV at
30~\AA\ and more than 1~eV on cells of 20~\AA, and should not be used as references or
regression values. For 2D the large-cell code should be run $\Gamma$-only
(\texttt{KPOINTS\_W 1 1 1}), which activates the head-and-wing correction; the same holds for
3D cells (\texttt{PERIODIC XYZ}), where $V(0)$ is an Ewald sum with the $\bG=0$ term omitted.

\subsection{Gaussian form factors of the auxiliary charges are not needed}
\label{sec:formfactor}

The 2D correction treats the auxiliary Gaussians as point charges at their centres, the limit
$q_\mathrm{max}\sigma_P\ll1$ of Eq.~\eqref{eq:Zl0}; the Gaussian form factors
$\ee^{-q^2/4\alpha}$ are neglected. With the most diffuse $s$ exponent of the aug-SZV RI basis
($\alpha = 0.199$) the head kernel changes by 7\,\% on the $3\times3$ cell, 1.8\,\% at
$6\times6$, 1.1\,\% at $8\times8$ and 0.4\,\% at $14\times14$. On a head contribution of
$\sim30$~meV to the gap this is below 1~meV on production cells, and the form factors are not
implemented. In 3D the question does not arise: the correction is applied in the orbital basis,
Eq.~\eqref{eq:What}, and involves no auxiliary charges.

\subsection{Remaining approximations}
\label{sec:disc_approx}
\label{sec:approx}

Four approximations of the model remain. (i) The body of $W$ is taken at $\Gamma$ only, so its
dispersion across the small Brillouin zone is not corrected; this error decreases with the cell
size, and the cells for which the method is intended have a flat body. (ii) In 2D, the
constant-function head carries no $\ee^{-q|z-z'|}$ form factor for the finite thickness of the
layer: a state that extends several \AA\ into the vacuum couples to the charge mode with a
factor $\ee^{-q|z|}$ that is $0.7$ at $q=0.1$~\AA$^{-1}$ and $z=4$~\AA. This was the first
candidate for the AlN difference of Sec.~\ref{sec:res2d}, but it does not survive the cell
test: the missing factor is first order in $q$ and its effect would shrink with the radius of
the small Brillouin zone, by about one half from the $9\times9$ to the $12\times12$ cell,
whereas the gap moves by 4~meV; and the exact-RI tensor driver reproduces the RI-RS value, so no
fit error is involved either. Our integral runs over the whole small Brillouin zone
($q\le\pi/9a$ for a $9\times9$ cell) while the PR applies its model only to the cells adjacent
to $\Gamma$ ($q\le\pi/36a$); which of the two $\Gamma$-point treatments is closer to the
full-$k$ limit for a diffuse state is not settled by the present data. The difference is not a
property of the box of the lattice sum: the subtracted constant
$D_\mathrm{box}$ grows with the box while the finite remainder of the charge channel does not
(h-BN, $8\times8$: $s=4$ and $s=5$ boxes differ by 1~meV in the gap). (iii) In 3D, the
small-$q$ expansion of the pair phase, Eq.~\eqref{eq:pair_phi0}, holds for $|\bq\cdot\mathbf
d_{nm}|<1$; nearly degenerate pairs with a finite coupling violate it, and their dipoles are
clipped at $1/q_{\max}$ (484 of $512^2$ pairs in the 128-atom diamond cell). (iv) The
orbital-diagonal application neglects the off-diagonal elements $\Delta\Sigma^c_{nm}$,
$n\neq m$, which only matter if the correction were allowed to mix quasiparticle states.

%======================================================================================
\section{Conclusions}
\label{sec:conclusions}
%======================================================================================

$\Gamma$-only RI-RS $GW$ of periodic large cells is made well defined by separating the charge
channel of the auxiliary basis. In 2D slabs its divergent lattice sum is removed from $V(0)$
as a single rank-one constant, $D_\mathrm{box}\,cc^\Tr$, computed exactly as a boundary integral
of the image box; in 3D bulk cells $V(0)$ is the Ewald sum with the $\bG=0$ term omitted. The
charge channel is put back analytically through the normalized constant function, whose
screened head and wings follow in closed form from the velocity-gauge transition dipoles and the
Gaussian dielectric body and are integrated over the whole small Brillouin zone of the
supercell; nothing is subtracted. In 2D the correction is applied in the auxiliary basis, at
negligible cost. In 3D, where the head is eV-sized and an RI-space application is amplified by
the fit, it is applied to $\Sigma^c$ in the orbital basis through a single orbital-pair kernel
that contains the head, the head-dipole and the wing terms. The implementation in \cpk{} is
activated automatically by \texttt{KPOINTS\_W 1 1 1} for \texttt{PERIODIC XY} and
\texttt{PERIODIC XYZ}, and chooses replicated or distributed dense algebra from a memory
estimate. The earlier $k$-mesh large-cell scheme carries a spurious charge leak of the fitted
pair densities and should not be used as a reference for 2D systems.
Numerically, with one master build on both sides, the slab correction reproduces the
converged full-$k$ RI-RS $GW$ of PR~\#6447 to 20--75~meV for h-BN (two bases), SiC and
h-GaN, with a 0.28~eV deviation for h-AlN that the exact-RI tensor driver reproduces to 1~meV,
so it is a property of the $\Gamma$-point models and not of the RI-RS fit. For diamond the
$\Gamma$-only RI-RS result of the 128-atom cell, exchange, correlation and correction, is
identical to the exact-RI tensor result to the meV, and the dipole terms of the bulk kernel
scale as $1/\Omega$ as derived; but the gap is still rising by 0.13~eV from 128 to 250 atoms
and lies 0.5~eV above the converged full-$k$ value there, entirely in $\Sigma^c$. The bulk
$\Gamma$-point model is therefore validated as an implementation and not yet as a reference:
which of the two $\Gamma$-point treatments is closer to the full-$k$ limit remains open, and
the standard tensor $k$-point reference proved too costly to decide it. The correction costs
about 1~\% of a 128-atom run.

\begin{acknowledgments}
The full-$k$ reference calculations use the implementation of J.~Wilhelm (cp2k PR~\#6447).
Computations were carried out on the Noctua~2 cluster of the Paderborn Center for Parallel
Computing (PC$^2$).
\end{acknowledgments}

\begin{thebibliography}{2}
\bibitem{wilhelm_fullk}
J.~Wilhelm, \emph{Periodic real-space RI for full-$k$ $GW$ band structures}, manuscript in
preparation.
\bibitem{pr6447}
\cpk{} pull request \#6447, \url{https://github.com/cp2k/cp2k/pull/6447}.
\end{thebibliography}

\end{document}
